% Options for packages loaded elsewhere
\PassOptionsToPackage{unicode}{hyperref}
\PassOptionsToPackage{hyphens}{url}
\PassOptionsToPackage{dvipsnames,svgnames,x11names}{xcolor}
%
\documentclass[
  11pt,
  a4paper,
]{scrbook}

\usepackage{amsmath,amssymb}
\usepackage{iftex}
\ifPDFTeX
  \usepackage[T1]{fontenc}
  \usepackage[utf8]{inputenc}
  \usepackage{textcomp} % provide euro and other symbols
\else % if luatex or xetex
  \usepackage{unicode-math}
  \defaultfontfeatures{Scale=MatchLowercase}
  \defaultfontfeatures[\rmfamily]{Ligatures=TeX,Scale=1}
\fi
\usepackage{lmodern}
\ifPDFTeX\else  
    % xetex/luatex font selection
\fi
% Use upquote if available, for straight quotes in verbatim environments
\IfFileExists{upquote.sty}{\usepackage{upquote}}{}
\IfFileExists{microtype.sty}{% use microtype if available
  \usepackage[]{microtype}
  \UseMicrotypeSet[protrusion]{basicmath} % disable protrusion for tt fonts
}{}
\makeatletter
\@ifundefined{KOMAClassName}{% if non-KOMA class
  \IfFileExists{parskip.sty}{%
    \usepackage{parskip}
  }{% else
    \setlength{\parindent}{0pt}
    \setlength{\parskip}{6pt plus 2pt minus 1pt}}
}{% if KOMA class
  \KOMAoptions{parskip=half}}
\makeatother
\usepackage{xcolor}
\usepackage[margin=2.5cm]{geometry}
\setlength{\emergencystretch}{3em} % prevent overfull lines
\setcounter{secnumdepth}{5}
% Make \paragraph and \subparagraph free-standing
\ifx\paragraph\undefined\else
  \let\oldparagraph\paragraph
  \renewcommand{\paragraph}[1]{\oldparagraph{#1}\mbox{}}
\fi
\ifx\subparagraph\undefined\else
  \let\oldsubparagraph\subparagraph
  \renewcommand{\subparagraph}[1]{\oldsubparagraph{#1}\mbox{}}
\fi

\usepackage{color}
\usepackage{fancyvrb}
\newcommand{\VerbBar}{|}
\newcommand{\VERB}{\Verb[commandchars=\\\{\}]}
\DefineVerbatimEnvironment{Highlighting}{Verbatim}{commandchars=\\\{\}}
% Add ',fontsize=\small' for more characters per line
\usepackage{framed}
\definecolor{shadecolor}{RGB}{241,243,245}
\newenvironment{Shaded}{\begin{snugshade}}{\end{snugshade}}
\newcommand{\AlertTok}[1]{\textcolor[rgb]{0.68,0.00,0.00}{#1}}
\newcommand{\AnnotationTok}[1]{\textcolor[rgb]{0.37,0.37,0.37}{#1}}
\newcommand{\AttributeTok}[1]{\textcolor[rgb]{0.40,0.45,0.13}{#1}}
\newcommand{\BaseNTok}[1]{\textcolor[rgb]{0.68,0.00,0.00}{#1}}
\newcommand{\BuiltInTok}[1]{\textcolor[rgb]{0.00,0.23,0.31}{#1}}
\newcommand{\CharTok}[1]{\textcolor[rgb]{0.13,0.47,0.30}{#1}}
\newcommand{\CommentTok}[1]{\textcolor[rgb]{0.37,0.37,0.37}{#1}}
\newcommand{\CommentVarTok}[1]{\textcolor[rgb]{0.37,0.37,0.37}{\textit{#1}}}
\newcommand{\ConstantTok}[1]{\textcolor[rgb]{0.56,0.35,0.01}{#1}}
\newcommand{\ControlFlowTok}[1]{\textcolor[rgb]{0.00,0.23,0.31}{#1}}
\newcommand{\DataTypeTok}[1]{\textcolor[rgb]{0.68,0.00,0.00}{#1}}
\newcommand{\DecValTok}[1]{\textcolor[rgb]{0.68,0.00,0.00}{#1}}
\newcommand{\DocumentationTok}[1]{\textcolor[rgb]{0.37,0.37,0.37}{\textit{#1}}}
\newcommand{\ErrorTok}[1]{\textcolor[rgb]{0.68,0.00,0.00}{#1}}
\newcommand{\ExtensionTok}[1]{\textcolor[rgb]{0.00,0.23,0.31}{#1}}
\newcommand{\FloatTok}[1]{\textcolor[rgb]{0.68,0.00,0.00}{#1}}
\newcommand{\FunctionTok}[1]{\textcolor[rgb]{0.28,0.35,0.67}{#1}}
\newcommand{\ImportTok}[1]{\textcolor[rgb]{0.00,0.46,0.62}{#1}}
\newcommand{\InformationTok}[1]{\textcolor[rgb]{0.37,0.37,0.37}{#1}}
\newcommand{\KeywordTok}[1]{\textcolor[rgb]{0.00,0.23,0.31}{#1}}
\newcommand{\NormalTok}[1]{\textcolor[rgb]{0.00,0.23,0.31}{#1}}
\newcommand{\OperatorTok}[1]{\textcolor[rgb]{0.37,0.37,0.37}{#1}}
\newcommand{\OtherTok}[1]{\textcolor[rgb]{0.00,0.23,0.31}{#1}}
\newcommand{\PreprocessorTok}[1]{\textcolor[rgb]{0.68,0.00,0.00}{#1}}
\newcommand{\RegionMarkerTok}[1]{\textcolor[rgb]{0.00,0.23,0.31}{#1}}
\newcommand{\SpecialCharTok}[1]{\textcolor[rgb]{0.37,0.37,0.37}{#1}}
\newcommand{\SpecialStringTok}[1]{\textcolor[rgb]{0.13,0.47,0.30}{#1}}
\newcommand{\StringTok}[1]{\textcolor[rgb]{0.13,0.47,0.30}{#1}}
\newcommand{\VariableTok}[1]{\textcolor[rgb]{0.07,0.07,0.07}{#1}}
\newcommand{\VerbatimStringTok}[1]{\textcolor[rgb]{0.13,0.47,0.30}{#1}}
\newcommand{\WarningTok}[1]{\textcolor[rgb]{0.37,0.37,0.37}{\textit{#1}}}

\providecommand{\tightlist}{%
  \setlength{\itemsep}{0pt}\setlength{\parskip}{0pt}}\usepackage{longtable,booktabs,array}
\usepackage{calc} % for calculating minipage widths
% Correct order of tables after \paragraph or \subparagraph
\usepackage{etoolbox}
\makeatletter
\patchcmd\longtable{\par}{\if@noskipsec\mbox{}\fi\par}{}{}
\makeatother
% Allow footnotes in longtable head/foot
\IfFileExists{footnotehyper.sty}{\usepackage{footnotehyper}}{\usepackage{footnote}}
\makesavenoteenv{longtable}
\usepackage{graphicx}
\makeatletter
\def\maxwidth{\ifdim\Gin@nat@width>\linewidth\linewidth\else\Gin@nat@width\fi}
\def\maxheight{\ifdim\Gin@nat@height>\textheight\textheight\else\Gin@nat@height\fi}
\makeatother
% Scale images if necessary, so that they will not overflow the page
% margins by default, and it is still possible to overwrite the defaults
% using explicit options in \includegraphics[width, height, ...]{}
\setkeys{Gin}{width=\maxwidth,height=\maxheight,keepaspectratio}
% Set default figure placement to htbp
\makeatletter
\def\fps@figure{htbp}
\makeatother

\makeatletter
\@ifpackageloaded{tcolorbox}{}{\usepackage[skins,breakable]{tcolorbox}}
\@ifpackageloaded{fontawesome5}{}{\usepackage{fontawesome5}}
\definecolor{quarto-callout-color}{HTML}{909090}
\definecolor{quarto-callout-note-color}{HTML}{0758E5}
\definecolor{quarto-callout-important-color}{HTML}{CC1914}
\definecolor{quarto-callout-warning-color}{HTML}{EB9113}
\definecolor{quarto-callout-tip-color}{HTML}{00A047}
\definecolor{quarto-callout-caution-color}{HTML}{FC5300}
\definecolor{quarto-callout-color-frame}{HTML}{acacac}
\definecolor{quarto-callout-note-color-frame}{HTML}{4582ec}
\definecolor{quarto-callout-important-color-frame}{HTML}{d9534f}
\definecolor{quarto-callout-warning-color-frame}{HTML}{f0ad4e}
\definecolor{quarto-callout-tip-color-frame}{HTML}{02b875}
\definecolor{quarto-callout-caution-color-frame}{HTML}{fd7e14}
\makeatother
\makeatletter
\@ifpackageloaded{bookmark}{}{\usepackage{bookmark}}
\makeatother
\makeatletter
\@ifpackageloaded{caption}{}{\usepackage{caption}}
\AtBeginDocument{%
\ifdefined\contentsname
  \renewcommand*\contentsname{Table of contents}
\else
  \newcommand\contentsname{Table of contents}
\fi
\ifdefined\listfigurename
  \renewcommand*\listfigurename{List of Figures}
\else
  \newcommand\listfigurename{List of Figures}
\fi
\ifdefined\listtablename
  \renewcommand*\listtablename{List of Tables}
\else
  \newcommand\listtablename{List of Tables}
\fi
\ifdefined\figurename
  \renewcommand*\figurename{Figure}
\else
  \newcommand\figurename{Figure}
\fi
\ifdefined\tablename
  \renewcommand*\tablename{Table}
\else
  \newcommand\tablename{Table}
\fi
}
\@ifpackageloaded{float}{}{\usepackage{float}}
\floatstyle{ruled}
\@ifundefined{c@chapter}{\newfloat{codelisting}{h}{lop}}{\newfloat{codelisting}{h}{lop}[chapter]}
\floatname{codelisting}{Listing}
\newcommand*\listoflistings{\listof{codelisting}{List of Listings}}
\makeatother
\makeatletter
\makeatother
\makeatletter
\@ifpackageloaded{caption}{}{\usepackage{caption}}
\@ifpackageloaded{subcaption}{}{\usepackage{subcaption}}
\makeatother
\ifLuaTeX
  \usepackage{selnolig}  % disable illegal ligatures
\fi
\usepackage{bookmark}

\IfFileExists{xurl.sty}{\usepackage{xurl}}{} % add URL line breaks if available
\urlstyle{same} % disable monospaced font for URLs
\hypersetup{
  pdftitle={MA Thesis Guide},
  pdfauthor={Philipp Masur},
  colorlinks=true,
  linkcolor={003D6B},
  filecolor={Maroon},
  citecolor={Blue},
  urlcolor={00A6D6},
  pdfcreator={LaTeX via pandoc}}

\title{MA Thesis Guide}
\usepackage{etoolbox}
\makeatletter
\providecommand{\subtitle}[1]{% add subtitle to \maketitle
  \apptocmd{\@title}{\par {\large #1 \par}}{}{}
}
\makeatother
\subtitle{Communication Science · Vrije Universiteit Amsterdam}
\author{Philipp Masur}
\date{March 18, 2026}

\begin{document}
\frontmatter
\maketitle

\renewcommand*\contentsname{Table of contents}
{
\hypersetup{linkcolor=}
\setcounter{tocdepth}{2}
\tableofcontents
}
\mainmatter
\bookmarksetup{startatroot}

\chapter*{Welcome}\label{welcome}
\addcontentsline{toc}{chapter}{Welcome}

\markboth{Welcome}{Welcome}

Welcome to the \textbf{MA Thesis Guide for Communication Science} at
Vrije Universiteit Amsterdam. This guide is designed to support you
throughout your master's thesis project --- from choosing a topic to
submitting your final version.

The guide brings together information from several official documents
(the thesis manual, writing guidelines, and departmental AI policy) in
one accessible, searchable format. Use it alongside your supervisor's
feedback and the resources on Canvas.

\section*{How to use this guide}\label{how-to-use-this-guide}
\addcontentsline{toc}{section}{How to use this guide}

\markright{How to use this guide}

This guide is organized into five parts:

\begin{enumerate}
\def\labelenumi{\arabic{enumi}.}
\tightlist
\item
  \textbf{Getting Started} --- course details, the thesis trajectory,
  key deadlines, and how supervision works
\item
  \textbf{The Research Process} --- from research question to ethics and
  data management
\item
  \textbf{Writing Your Thesis} --- section-by-section writing
  requirements and tips
\item
  \textbf{Academic Practices} --- literature search, APA citation, open
  science, and AI tools
\item
  \textbf{Submission \& Assessment} --- grading, resits, and graduation
\end{enumerate}

Use the sidebar to navigate, the search bar to find specific topics, or
download the full guide as a PDF or EPUB using the buttons above.

\begin{tcolorbox}[enhanced jigsaw, bottomrule=.15mm, opacityback=0, toprule=.15mm, colframe=quarto-callout-tip-color-frame, left=2mm, rightrule=.15mm, title=\textcolor{quarto-callout-tip-color}{\faLightbulb}\hspace{0.5em}{New to the thesis process?}, colbacktitle=quarto-callout-tip-color!10!white, bottomtitle=1mm, colback=white, breakable, toptitle=1mm, titlerule=0mm, leftrule=.75mm, arc=.35mm, opacitybacktitle=0.6, coltitle=black]

Start with \href{getting-started.qmd}{Getting Started} for an overview
of the course, then work through \href{trajectory.qmd}{The Thesis
Trajectory} to understand what is expected at each stage.

\end{tcolorbox}

\begin{tcolorbox}[enhanced jigsaw, bottomrule=.15mm, opacityback=0, toprule=.15mm, colframe=quarto-callout-note-color-frame, left=2mm, rightrule=.15mm, title=\textcolor{quarto-callout-note-color}{\faInfo}\hspace{0.5em}{Official documents}, colbacktitle=quarto-callout-note-color!10!white, bottomtitle=1mm, colback=white, breakable, toptitle=1mm, titlerule=0mm, leftrule=.75mm, arc=.35mm, opacitybacktitle=0.6, coltitle=black]

This guide synthesizes and elaborates on the following official
departmental documents:

\begin{itemize}
\tightlist
\item
  \emph{Manual for the MA thesis Communication Science 2025--26}
\item
  \emph{Guideline Masterthesis Communication Science}
\item
  \emph{Use of Generative AI Tools in Your Ba/Ma Thesis 2025--26}
\end{itemize}

In case of any discrepancy, the official documents on Canvas take
precedence.

\end{tcolorbox}

\section*{About this guide}\label{about-this-guide}
\addcontentsline{toc}{section}{About this guide}

\markright{About this guide}

This guide was created by \textbf{Philipp Masur} (Department of
Communication Science, VU Amsterdam). It is updated annually to reflect
current departmental policy and best practices.

If you notice an error or want to suggest an improvement, use the ``Edit
this page'' or ``Report an issue'' links in the right margin of any
page.

\part{Getting Started}

\chapter{Getting Started}\label{getting-started-1}

\section{Course information}\label{course-information}

The MA Thesis in Communication Science is a substantial independent
research project that forms the capstone of your master's degree. Here
are the key course facts:

\begin{longtable}[]{@{}
  >{\raggedright\arraybackslash}p{(\columnwidth - 2\tabcolsep) * \real{0.5000}}
  >{\raggedright\arraybackslash}p{(\columnwidth - 2\tabcolsep) * \real{0.5000}}@{}}
\toprule\noalign{}
\endhead
\bottomrule\noalign{}
\endlastfoot
\textbf{Course code} & S\_MTcw \\
\textbf{Credits} & 18 EC (≈ 504 hours) \\
\textbf{Course level} & 600 \\
\textbf{Periods} & 4--6 (February track) or 1--2 (September track) \\
\textbf{Language} & Dutch or English \\
\textbf{Faculty} & Faculty of Social Sciences \\
\textbf{Coordinator} & dr. Camiel Beukeboom \\
\textbf{Teaching format} & Individual and small-group meetings with
supervisor, plus centralized support and student peer review \\
\textbf{Assessment} & Evaluated by supervisor (first reviewer) and a
second reviewer using the assessment form (see Canvas) \\
\end{longtable}

\section{Entry requirements}\label{entry-requirements}

Before you can continue with the full MA thesis trajectory, you must
meet two requirements:

\begin{enumerate}
\def\labelenumi{\arabic{enumi}.}
\tightlist
\item
  A minimum of \textbf{12 EC} in total from preceding MA courses,
  including successful completion of \textbf{Research Methods for
  Communication Science}
\item
  An \textbf{approved Research Proposal} after the first phase (see
  \href{trajectory.qmd}{The Thesis Trajectory})
\end{enumerate}

\section{Start dates: February or September
track}\label{start-dates-february-or-september-track}

The thesis course formally starts in \textbf{early February} (Period 4).
However, it is also possible to start after the summer, in \textbf{early
September}.

\textbf{Why choose the September track?}

\begin{itemize}
\tightlist
\item
  You are doing an internship or other study activity that does not fit
  alongside the thesis
\item
  You did not pass all courses from Period 1 or 2
\end{itemize}

If you want to start in September, discuss this with potential
supervisors during the thesis market in January. Note that only a few
supervisors supervise theses in the September track. You must still
participate in the January supervisor assignment procedure.

\section{Registration}\label{registration}

Preparation begins in \textbf{December} with an information lecture
about procedures, deadlines, and available supervisor topics.
\textbf{Register for the course before December} to receive information
via Canvas.

\begin{tcolorbox}[enhanced jigsaw, bottomrule=.15mm, opacityback=0, toprule=.15mm, colframe=quarto-callout-important-color-frame, left=2mm, rightrule=.15mm, title=\textcolor{quarto-callout-important-color}{\faExclamation}\hspace{0.5em}{Don't miss the information lecture}, colbacktitle=quarto-callout-important-color!10!white, bottomtitle=1mm, colback=white, breakable, toptitle=1mm, titlerule=0mm, leftrule=.75mm, arc=.35mm, opacitybacktitle=0.6, coltitle=black]

The information lecture (usually early December) is where you receive
the full list of supervisor topics and procedural information. Missing
it can mean missing important deadlines.

\end{tcolorbox}

\section{Choosing a topic and finding a
supervisor}\label{choosing-a-topic-and-finding-a-supervisor}

In \textbf{January}, a \textbf{thesis market} is organized where you can
walk around and talk informally with different supervisors about their
research themes.

After the market, you submit an \textbf{online questionnaire} with your
topic preferences (ranked 1--8). You can also upload your own research
proposal via the same form (but you should still indicate topic
preferences in case your own proposal is not feasible).

Supervisors are generally assigned by the \textbf{end of January}.

\subsection{Can I propose my own
topic?}\label{can-i-propose-my-own-topic}

Yes. However, not all proposals can be accommodated --- it must be
judged as interesting and feasible, and a supervisor must be available.
A good own proposal should contain at minimum:

\begin{itemize}
\tightlist
\item
  A research question
\item
  A brief theoretical foundation
\item
  The intended type of research (survey, experiment, content analysis,
  etc.)
\end{itemize}

\subsection{Can I write my thesis in
English?}\label{can-i-write-my-thesis-in-english}

Yes. You can write in Dutch or English, in consultation with your
supervisor. Note that some supervisors only accept theses in English. If
you wish to write in any other language, the Examination Board must
grant permission.

\chapter{The Thesis Trajectory}\label{the-thesis-trajectory}

The MA thesis project unfolds in three main phases: \textbf{preparation
(P4)}, \textbf{research and writing (P5)}, and \textbf{finalizing and
submission (P6)}. This chapter explains what is expected at each stage.

\section{Phase 1 (P4): Preparation}\label{phase-1-p4-preparation}

Period 4 runs from \textbf{February to late March}. During this period,
the thesis course runs in parallel with your two elective courses ---
plan to spend approximately \textbf{4 hours per week} on your thesis in
this phase.

Your main tasks in P4:

\begin{itemize}
\tightlist
\item
  Read relevant literature in your chosen area
\item
  Delineate your individual research focus
\item
  Develop a draft research question
\item
  Meet with your supervisor (and thesis group) to get feedback on your
  ideas
\end{itemize}

\begin{tcolorbox}[enhanced jigsaw, bottomrule=.15mm, opacityback=0, toprule=.15mm, colframe=quarto-callout-tip-color-frame, left=2mm, rightrule=.15mm, title=\textcolor{quarto-callout-tip-color}{\faLightbulb}\hspace{0.5em}{Start writing immediately}, colbacktitle=quarto-callout-tip-color!10!white, bottomtitle=1mm, colback=white, breakable, toptitle=1mm, titlerule=0mm, leftrule=.75mm, arc=.35mm, opacitybacktitle=0.6, coltitle=black]

As soon as you start reading, write down your ideas and summaries. Keep
an overview of sources (e.g., in Zotero) and note the relevant aspects
of each paper. This prevents you from getting lost in the literature and
speeds up writing later.

\end{tcolorbox}

\subsection{The MA Thesis Plan (end of
P4)}\label{the-ma-thesis-plan-end-of-p4}

By the \textbf{end of Period 4} (see \href{timeline.qmd}{Timeline \&
Deadlines}), you submit a \textbf{Thesis Plan} to your supervisor via
email. This is a \textbf{minimum two-page document} that outlines your
proposed project. The further you have developed it, the better. It
should include at minimum:

\begin{enumerate}
\def\labelenumi{\arabic{enumi}.}
\tightlist
\item
  A first description of the relevant literature and core theory, and a
  research ``gap'' that justifies your focus (theoretical and societal
  relevance)
\item
  Your research question
\item
  Related hypotheses or sub-questions
\item
  A description of your intended research design (e.g., \emph{``a survey
  exploring the relationship between X and Y''}, or \emph{``a 2×2
  between-subjects experiment with X as IV and Y as DV''})
\end{enumerate}

Think of the Thesis Plan as a blueprint for your Research Proposal: (1)
and (2) become your Introduction, (3) goes into the Theory section, and
(4) goes into the Method section.

\section{Phase 2 (P5): Research and
Writing}\label{phase-2-p5-research-and-writing}

Period 5 begins in April. You now work \textbf{full time} on your
thesis.

\subsection{The Research Proposal (early
P5)}\label{the-research-proposal-early-p5}

In the \textbf{first two weeks of P5}, you intensively develop your
\textbf{Research Proposal}, building on the Thesis Plan and
incorporating feedback from your supervisor. The Research Proposal must
contain:

\begin{enumerate}
\def\labelenumi{\arabic{enumi}.}
\tightlist
\item
  \textbf{Introduction} --- including a clearly formulated Research
  Question
\item
  \textbf{Theory} --- theoretical grounds and argumentation for the
  relationships you expect
\item
  \textbf{Final hypotheses} (or sub-questions for qualitative research)
\item
  A \textbf{conceptual model} (for quantitative research)
\item
  \textbf{Method} --- research design, (in)dependent and control
  variables, ideas on procedure and operationalization (a first draft of
  the questionnaire, codebook, etc. can go in the appendix)
\item
  \textbf{Analysis plan} --- for each hypothesis/RQ, specify exactly
  which statistical test on which variables you will use
\end{enumerate}

\begin{tcolorbox}[enhanced jigsaw, bottomrule=.15mm, opacityback=0, toprule=.15mm, colframe=quarto-callout-important-color-frame, left=2mm, rightrule=.15mm, title=\textcolor{quarto-callout-important-color}{\faExclamation}\hspace{0.5em}{The Research Proposal is a gate}, colbacktitle=quarto-callout-important-color!10!white, bottomtitle=1mm, colback=white, breakable, toptitle=1mm, titlerule=0mm, leftrule=.75mm, arc=.35mm, opacitybacktitle=0.6, coltitle=black]

Your supervisor evaluates the proposal as \emph{sufficient} or
\emph{insufficient}. An \textbf{approved Research Proposal is required
to continue} the thesis trajectory. If it is rejected, you have up to
two resits (within 10 working days each). Failing both resits means you
must stop and restart in the following academic year.

\end{tcolorbox}

Upload your Research Proposal on Canvas \textbf{and} email it to your
supervisor before the deadline. Give peer feedback to fellow students
first, and process their feedback before submitting.

\subsection{Complete Methodology}\label{complete-methodology}

As soon as possible after the Research Proposal is approved, submit your
\textbf{Complete Methodology} to your supervisor. This is the final,
detailed operationalization of your method --- the exact content of your
questionnaire, experimental materials, codebook, Qualtrics setup, etc.
This document usually becomes an appendix in your final thesis.

\subsection{Ethical Review}\label{ethical-review}

Before you start collecting data, you must complete the \textbf{ethical
self-check} via the Research Ethics Review Committee (RERC). See
\href{ethics.qmd}{Ethics \& Data Management} for details.

\subsection{Data Collection and
Analysis}\label{data-collection-and-analysis}

Once your Research Proposal and Complete Methodology are approved and
the ethical self-check is done:

\begin{itemize}
\tightlist
\item
  Start \textbf{data collection}
\item
  Write the \textbf{Method chapter} (in past tense)
\item
  Conduct analyses and write the \textbf{Results chapter}
\end{itemize}

\begin{tcolorbox}[enhanced jigsaw, bottomrule=.15mm, opacityback=0, toprule=.15mm, colframe=quarto-callout-tip-color-frame, left=2mm, rightrule=.15mm, title=\textcolor{quarto-callout-tip-color}{\faLightbulb}\hspace{0.5em}{Prepare analyses before data comes in}, colbacktitle=quarto-callout-tip-color!10!white, bottomtitle=1mm, colback=white, breakable, toptitle=1mm, titlerule=0mm, leftrule=.75mm, arc=.35mm, opacitybacktitle=0.6, coltitle=black]

You can write your analysis syntax (R or SPSS) and structure your
results section before the final data is collected. This saves a lot of
time.

\end{tcolorbox}

\section{Phase 3 (P6): Finalizing and
Submission}\label{phase-3-p6-finalizing-and-submission}

Period 6 is the final stretch. After receiving feedback on your Method
and Results chapters, you:

\begin{itemize}
\tightlist
\item
  Write the \textbf{Conclusion and Discussion} chapter
\item
  Polish Introduction, Theory, Method, and Results
\item
  Process all remaining feedback
\item
  Submit the \textbf{complete concept version} to your supervisor for
  final feedback

  \begin{itemize}
  \tightlist
  \item
    This round focuses on: (1) overall coherence, (2) the
    discussion/conclusion section, and (3) any specific questions you
    have
  \end{itemize}
\item
  After processing all feedback, submit the \textbf{final version}
\end{itemize}

\begin{tcolorbox}[enhanced jigsaw, bottomrule=.15mm, opacityback=0, toprule=.15mm, colframe=quarto-callout-note-color-frame, left=2mm, rightrule=.15mm, title=\textcolor{quarto-callout-note-color}{\faInfo}\hspace{0.5em}{Quality matters --- for you and your supervisor}, colbacktitle=quarto-callout-note-color!10!white, bottomtitle=1mm, colback=white, breakable, toptitle=1mm, titlerule=0mm, leftrule=.75mm, arc=.35mm, opacitybacktitle=0.6, coltitle=black]

Submit only high-quality draft versions to your supervisor, meaning:
versions you have already peer-reviewed with fellow students and
improved. Feedback can only be useful and focused if the text is already
in good shape.

\end{tcolorbox}

\section{Peer feedback}\label{peer-feedback}

Throughout the trajectory, you are expected to give and receive
\textbf{peer feedback} with fellow students in your thesis group. Peer
feedback is required before submitting the following to your supervisor:

\begin{itemize}
\tightlist
\item
  The Research Proposal
\item
  The Method and Results chapters
\item
  The complete concept version
\end{itemize}

Use the official \textbf{Guideline Masterthesis Communication Science}
(Canvas) as your peer feedback checklist.

\chapter{Timeline \& Key Deadlines}\label{timeline-key-deadlines}

\section{Monthly overview}\label{monthly-overview}

The table below shows the typical tasks and tips for each month in the
\textbf{February track} (P4--P6). The same structure applies to the
September track, shifted by approximately seven months.

\begin{longtable}[]{@{}
  >{\raggedright\arraybackslash}p{(\columnwidth - 4\tabcolsep) * \real{0.3500}}
  >{\raggedright\arraybackslash}p{(\columnwidth - 4\tabcolsep) * \real{0.3500}}
  >{\raggedright\arraybackslash}p{(\columnwidth - 4\tabcolsep) * \real{0.3000}}@{}}
\toprule\noalign{}
\begin{minipage}[b]{\linewidth}\raggedright
Month
\end{minipage} & \begin{minipage}[b]{\linewidth}\raggedright
Tasks
\end{minipage} & \begin{minipage}[b]{\linewidth}\raggedright
Tips
\end{minipage} \\
\midrule\noalign{}
\endhead
\bottomrule\noalign{}
\endlastfoot
\textbf{December} & Thesis information lecture & Think about what you
want to study. A preliminary list of supervisor topics is published
before Christmas. \\
\textbf{January} & Thesis market --- meet supervisors, discuss topics,
submit preferences questionnaire & Come prepared and with an open mind.
Talk to a range of supervisors. \\
\textbf{February (P4 starts)} & Read relevant literature; delineate your
individual focus; develop a draft research question & Start writing down
your ideas immediately. Use Zotero to keep track of sources. \\
\textbf{March} & Determine your research gap; formulate your research
question; think about your method. Submit \textbf{MA Thesis Plan}
(2-pager) to supervisor by end of P4. & Everything you do in P4 is
preparation for a flying start in P5. \\
\textbf{April (P5 starts)} & First two weeks: intensively develop your
\textbf{Research Proposal} (intro, theory, method, analysis plan).
Process peer feedback before submitting. & Use the knowledge clips and
support resources on Canvas for writing and hypothesis formulation. \\
\textbf{April} & Work out your \textbf{Complete Methodology}. Once
Research Proposal + analysis plan + methodology are approved: do the
ethical self-check, and optionally submit a pre-registration. Then start
data collection. & \\
\textbf{Late April} & Start data collection. Write method chapter;
conduct and report analyses in results chapter. & Keep a clean SPSS/R
syntax from the start. You can prepare your analysis syntax before data
is in. \\
\textbf{Early May} & Process peer feedback on method and results
chapters before submitting to supervisor. & Relate results to hypotheses
and theory; seek additional literature if needed. \\
\textbf{May} & Write and polish theory, method, results; write
discussion and conclusion. Process peer feedback on complete concept
version before submitting. & Writing a good thesis always requires many
revisions. \\
\textbf{Early June} & Submit complete concept version to supervisor.
Process feedback and refine. & Agree with your supervisor in advance
when to submit --- this is a busy time for everyone. \\
\textbf{End of June} & Submit \textbf{final version} & See deadline
table below. \\
\end{longtable}

\section{Deadline tables (2025--26)}\label{deadline-tables-202526}

\subsection{February track}\label{february-track}

\begin{longtable}[]{@{}ll@{}}
\toprule\noalign{}
Event & Date \\
\midrule\noalign{}
\endhead
\bottomrule\noalign{}
\endlastfoot
Information lecture & Tuesday, 2 December 2025 \\
Thesis market & Tuesday, 13 January 2026 \\
MA Thesis Plan deadline & Tuesday, 31 March 2026 \\
Research Proposal deadline & Tuesday, 14 April 2026 \\
\textbf{1st submission deadline} & \textbf{Friday, 26 June 2026,
16:00} \\
Resit deadline & Friday, 24 July 2026, 17:00 \\
2nd resit (if allowed) & Friday, 20 November 2026 \\
\end{longtable}

\subsection{September track}\label{september-track}

\begin{longtable}[]{@{}ll@{}}
\toprule\noalign{}
Event & Date \\
\midrule\noalign{}
\endhead
\bottomrule\noalign{}
\endlastfoot
Supervisor assignment & January 2026 \\
Start of trajectory & Beginning of September 2026 \\
MA Thesis Plan deadline & Monday, 19 October 2026 \\
Research Proposal deadline & Monday, 2 November 2026 \\
\textbf{1st submission deadline} & \textbf{Friday, 15 January 2027,
17:00} \\
Resit deadline & Friday, 19 February 2027, 17:00 \\
2nd resit (if allowed) & Friday, 18 June 2027 \\
\end{longtable}

\begin{tcolorbox}[enhanced jigsaw, bottomrule=.15mm, opacityback=0, toprule=.15mm, colframe=quarto-callout-important-color-frame, left=2mm, rightrule=.15mm, title=\textcolor{quarto-callout-important-color}{\faExclamation}\hspace{0.5em}{Missing the first deadline = missing an exam opportunity}, colbacktitle=quarto-callout-important-color!10!white, bottomtitle=1mm, colback=white, breakable, toptitle=1mm, titlerule=0mm, leftrule=.75mm, arc=.35mm, opacitybacktitle=0.6, coltitle=black]

Not submitting at the first opportunity is treated the same as not
showing up for an exam. You are then entitled to one resit. If you fail
the resit as well, you have exhausted your attempts for the current
academic year.

\end{tcolorbox}

\section{Supervisor availability}\label{supervisor-availability}

Your supervisor is available for guidance during the regular thesis
period, \textbf{up to and including the first submission deadline}.

\begin{itemize}
\tightlist
\item
  \textbf{During July and August}: supervisors are generally not
  available for feedback
\item
  \textbf{September--November resit period}: supervisors are only
  limitedly available
\end{itemize}

Plan accordingly --- do not rely on supervisor feedback during the
summer months.

\chapter{Supervision \& Peer Feedback}\label{supervision-peer-feedback}

\section{What to expect from your
supervisor}\label{what-to-expect-from-your-supervisor}

Your supervisor will guide you in:

\begin{itemize}
\tightlist
\item
  Formulating an interesting and well-grounded research question and
  hypotheses
\item
  Writing a coherent, well-structured, and theoretically sound thesis
\item
  Designing a proper methodology and analysis plan
\item
  Giving focused feedback on each chapter --- \textbf{once per chapter}
\end{itemize}

As a rule of thumb, your supervisor will provide feedback on:

\begin{enumerate}
\def\labelenumi{\arabic{enumi}.}
\tightlist
\item
  Your \textbf{introduction} (as part of the Thesis Plan feedback)
\item
  Your \textbf{theory chapter and research design} (as part of the
  Research Proposal feedback)
\item
  Your \textbf{complete methodology}
\item
  Your \textbf{method and results chapters}
\item
  The \textbf{complete concept version} (focus: coherence and
  discussion/conclusion)
\end{enumerate}

\section{What your supervisor expects from
you}\label{what-your-supervisor-expects-from-you}

A master's thesis student is expected to demonstrate the capacity to
\textbf{independently} set up, conduct, and report an empirical research
project. Your supervisor is a guide --- not a co-author.

This means:

\begin{itemize}
\tightlist
\item
  Show \textbf{initiative}: contact your supervisor with a clear agenda,
  not just general questions
\item
  Submit \textbf{high-quality drafts}: only submit versions that you
  (and your peers) cannot improve further on your own
\item
  \textbf{Do not ask your supervisor} about basic matters such as
  spell-checking, APA formatting, how to run R/SPSS tests, or things you
  have already been taught in your coursework
\end{itemize}

For technical analysis questions (how to run tests in R or SPSS, how to
report results), use:

\begin{itemize}
\tightlist
\item
  The statistics helpdesk
\item
  The \emph{R for Social Scientists} Canvas page
\item
  The departmental knowledge clips on Canvas
\item
  The Guideline Masterthesis Communication Science document
\end{itemize}

\begin{tcolorbox}[enhanced jigsaw, bottomrule=.15mm, opacityback=0, toprule=.15mm, colframe=quarto-callout-warning-color-frame, left=2mm, rightrule=.15mm, title=\textcolor{quarto-callout-warning-color}{\faExclamationTriangle}\hspace{0.5em}{Feedback is only useful if you've done your part}, colbacktitle=quarto-callout-warning-color!10!white, bottomtitle=1mm, colback=white, breakable, toptitle=1mm, titlerule=0mm, leftrule=.75mm, arc=.35mm, opacitybacktitle=0.6, coltitle=black]

Feedback is pointless on text you know has flaws. Submit drafts only
when you have already done everything you can to improve them ---
including peer feedback.

\end{tcolorbox}

\section{Thesis groups and peer
feedback}\label{thesis-groups-and-peer-feedback}

Most supervisors organize supervision in \textbf{thesis groups} of 2--6
students working on related topics. You will meet together especially in
P4. As projects become more individualized, groups may split into
smaller subgroups. Across the full trajectory, expect approximately
\textbf{6--8 meetings} with your supervisor.

\subsection{Peer feedback process}\label{peer-feedback-process}

Before submitting any draft to your supervisor, you must:

\begin{enumerate}
\def\labelenumi{\arabic{enumi}.}
\tightlist
\item
  Give and receive peer feedback with fellow students
\item
  Process the peer feedback and improve your draft
\item
  \textbf{Then} submit the improved version to your supervisor
\end{enumerate}

Peer feedback is required for:

\begin{itemize}
\tightlist
\item
  The \textbf{Research Proposal} (including introduction and theory)
\item
  The \textbf{Method and Results chapters}
\item
  The \textbf{complete concept version} of the thesis
\end{itemize}

Use the \textbf{Guideline Masterthesis Communication Science} (Canvas)
as your peer feedback checklist. As a peer reviewer, reflect on each
section in a few sentences and suggest concrete improvements.

\section{Formal regulations}\label{formal-regulations}

The formal regulations for the MA thesis are described in the
\emph{Regulation Ma thesis} document, available on VUnet. In practice,
supervisor and student can usually find a good working relationship
without needing to consult the formal rules. However, in cases of
disagreement, the regulations serve as the official reference point.

\section{Getting help for personal
circumstances}\label{getting-help-for-personal-circumstances}

If personal circumstances are affecting your progress:

\begin{itemize}
\tightlist
\item
  Contact the \textbf{study advisor} as early as possible:
  \href{https://vu.nl/en/student/contact-student-guidance-and-support/academic-advisors-fss}{vu.nl/en/student/contact-student-guidance-and-support/academic-advisors-fss}
\item
  For delays due to illness, psychological issues, or other
  circumstances, you may be eligible for financial support from the
  \textbf{Profileringsfonds}:
  \href{https://vu.nl/en/student/finances/profile-fund}{vu.nl/en/student/finances/profile-fund}
\end{itemize}

\begin{tcolorbox}[enhanced jigsaw, bottomrule=.15mm, opacityback=0, toprule=.15mm, colframe=quarto-callout-tip-color-frame, left=2mm, rightrule=.15mm, title=\textcolor{quarto-callout-tip-color}{\faLightbulb}\hspace{0.5em}{Report delays early}, colbacktitle=quarto-callout-tip-color!10!white, bottomtitle=1mm, colback=white, breakable, toptitle=1mm, titlerule=0mm, leftrule=.75mm, arc=.35mm, opacitybacktitle=0.6, coltitle=black]

The sooner you contact your study advisor about a delay, the more
options are available to you. Waiting too long limits what support can
be offered.

\end{tcolorbox}

\chapter{Writing the Thesis Plan}\label{writing-the-thesis-plan}

The \textbf{MA Thesis Plan} (sometimes called the Exposé) is the first
formal deliverable of your thesis project. It is due at the end of
Period 4 and should be a minimum of two pages. Think of it as a working
blueprint: it does not need to be polished or complete, but it should
show that you have done serious preparatory thinking.

\section{What the Thesis Plan needs to
contain}\label{what-the-thesis-plan-needs-to-contain}

\subsection{1. Literature overview and research
gap}\label{literature-overview-and-research-gap}

Begin by summarizing the state of knowledge in your topic area. You do
not need an exhaustive review --- focus on the core literature directly
relevant to your research focus. Then clearly articulate the
\textbf{research gap}: what is not yet known, understudied,
inconsistent, or unresolved?

A good research gap statement follows this pattern:

\begin{quote}
\emph{``While prior research has shown that X, relatively little
attention has been paid to Y. In particular, it remains unclear
whether/how/to what extent {[}gap{]}. This matters because
{[}theoretical/societal reason{]}.''}
\end{quote}

\textbf{Typical length in the plan:} 1--2 paragraphs.

\subsection{2. Research question}\label{research-question}

State your central research question explicitly. At this stage it may
still be somewhat provisional --- you can refine it in the Research
Proposal. A good research question is:

\begin{itemize}
\tightlist
\item
  Specific about the constructs involved
\item
  Empirically answerable
\item
  Clearly connected to the gap you identified
\end{itemize}

\textbf{Example:}

\begin{quote}
\emph{``To what extent does the framing of climate change news as a
political versus scientific issue affect readers' policy support, and
does this effect depend on prior environmental attitudes?''}
\end{quote}

\subsection{3. Hypotheses or
sub-questions}\label{hypotheses-or-sub-questions}

List your expected hypotheses (quantitative research) or sub-questions
(qualitative research). At this stage, these can be tentative. You
should be able to provide at least a brief theoretical rationale for
each.

For quantitative research, a hypothesis predicts a specific directional
or non-directional relationship:

\begin{quote}
\emph{``H1: Exposure to politically framed climate news will be
associated with lower policy support than scientifically framed news.''}

\emph{``H2: This effect will be moderated by prior environmental
attitudes, such that the framing effect is stronger among politically
non-engaged readers.''}
\end{quote}

\subsection{4. Research design}\label{research-design}

Describe how you plan to answer your research question. Be concrete ---
specify the type of research and the key elements. For example:

\begin{quote}
\emph{``I plan to conduct a 2 (news framing: political vs.~scientific) ×
2 (headline valence: positive vs.~negative) between-subjects online
experiment. Participants will be recruited via Prolific. The main
dependent variable is policy support, measured with an established
5-item scale.''}
\end{quote}

or

\begin{quote}
\emph{``I will use semi-structured interviews with 15--20 journalists at
Dutch national news outlets to explore how they perceive and enact
editorial independence norms in the context of audience metrics.''}
\end{quote}

\section{What makes a strong Thesis
Plan?}\label{what-makes-a-strong-thesis-plan}

\begin{longtable}[]{@{}
  >{\raggedright\arraybackslash}p{(\columnwidth - 2\tabcolsep) * \real{0.5714}}
  >{\raggedright\arraybackslash}p{(\columnwidth - 2\tabcolsep) * \real{0.4286}}@{}}
\toprule\noalign{}
\begin{minipage}[b]{\linewidth}\raggedright
Strong
\end{minipage} & \begin{minipage}[b]{\linewidth}\raggedright
Weak
\end{minipage} \\
\midrule\noalign{}
\endhead
\bottomrule\noalign{}
\endlastfoot
Research gap is clearly derived from the literature & Gap is vague or
asserted without evidence \\
RQ is specific and names constructs & RQ is too broad or not
researchable \\
Hypotheses follow logically from the theoretical argument & Hypotheses
are not justified \\
Design is clearly matched to the RQ & Design is described only in
general terms \\
Writing is clear and coherent & Writing is unfocused, ideas are
scattered \\
\end{longtable}

\section{Common mistakes}\label{common-mistakes}

\begin{itemize}
\tightlist
\item
  \textbf{Describing an area of interest rather than a gap.} ``Social
  media and mental health is an important topic'' is not a gap. What
  specifically is unknown?
\item
  \textbf{Overly broad RQs.} If your RQ could be a book, it's too broad.
  Narrow down to one specific relationship or phenomenon.
\item
  \textbf{Missing the connection between design and RQ.} Your design
  must be capable of answering your RQ. If your RQ asks about causal
  effects, a cross-sectional survey alone is insufficient.
\item
  \textbf{Writing the plan after the fact.} The plan is meant to guide
  your proposal, not just describe what you decided already. Use it to
  think through whether your project is feasible.
\end{itemize}

\section{Practical tips}\label{practical-tips}

\begin{tcolorbox}[enhanced jigsaw, bottomrule=.15mm, opacityback=0, toprule=.15mm, colframe=quarto-callout-tip-color-frame, left=2mm, rightrule=.15mm, title=\textcolor{quarto-callout-tip-color}{\faLightbulb}\hspace{0.5em}{Start earlier than you think}, colbacktitle=quarto-callout-tip-color!10!white, bottomtitle=1mm, colback=white, breakable, toptitle=1mm, titlerule=0mm, leftrule=.75mm, arc=.35mm, opacitybacktitle=0.6, coltitle=black]

The Thesis Plan is due at the end of P4, but the earlier you start
reading and writing, the better your plan will be --- and the easier
your Research Proposal will become.

\end{tcolorbox}

\begin{itemize}
\tightlist
\item
  \textbf{Use Zotero from the start}: track every paper you read with a
  brief note on its relevance to your project
\item
  \textbf{Write in parallel with reading}: don't wait until you've read
  everything to start writing
\item
  \textbf{Discuss your gap with your supervisor early}: your supervisor
  can quickly tell you if your gap has already been addressed in
  literature you haven't found yet
\item
  \textbf{Keep it focused}: two pages forces you to be precise ---
  that's the point
\end{itemize}

\section{Template structure}\label{template-structure}

A two-page Thesis Plan might be organized as follows:

\begin{verbatim}
Title (working title of the thesis)
Name | Student number | Supervisor | Date

1. Introduction and research gap (approx. ½ page)
   Brief overview of the topic; what is known; what is the gap

2. Research question (1–2 sentences)

3. Hypotheses / sub-questions (bulleted list with brief rationale)

4. Proposed research design (approx. ½ page)
   Type of study; participants/corpus; key variables/concepts; method

[Optional: preliminary reference list]
\end{verbatim}

\chapter{Writing the Research
Proposal}\label{writing-the-research-proposal}

The \textbf{Research Proposal} is the most important early milestone of
your thesis. It is due in the first two weeks of Period 5 and is
evaluated by your supervisor as \emph{sufficient} or
\emph{insufficient}. An approved proposal is required to continue --- it
is the green light to start collecting data.

The proposal is not the final thesis. It is a detailed plan that shows
your research design is sound, your theory is well-grounded, and you
know exactly what you are going to do and how. Your supervisor evaluates
it primarily on \textbf{feasibility} and the \textbf{fit between your
method, analysis plan, and research question}.

\section{Required sections}\label{required-sections}

\subsection{1. Introduction}\label{introduction}

The introduction frames your study and states your research question. It
follows the same structure as the introduction in the final thesis (see
\href{writing-intro-theory.qmd}{Introduction \& Theory Chapter}), but at
this stage it may not yet be in final form.

The introduction must contain: - A clear \textbf{motivation} for the
topic (why does this matter?) - An identification of the
\textbf{research gap} - A clear, explicitly stated \textbf{research
question} - A brief argument for the \textbf{societal and scientific
relevance}

\textbf{Length:} 1--2 pages.

\textbf{Common mistakes:} - Stating why the topic is interesting in
general, without identifying a specific gap - Formulating a research
question that is too broad, too vague, or not empirically answerable -
Relevance that is disconnected from the gap (e.g., claiming societal
relevance for a purely theoretical contribution)

\subsection{2. Theory}\label{theory}

The theory chapter reviews relevant literature and builds a theoretical
argument that culminates in your hypotheses. In the proposal, this
chapter should be substantially developed, though it need not be fully
polished.

Structure: 1. Define and explain your central constructs 2. Review prior
empirical work, organized thematically 3. Build theoretical arguments
linking constructs 4. Formulate hypotheses (or sub-questions) 5. Present
a \textbf{conceptual model} (for quantitative research)

\textbf{Length:} 4--6 pages in the proposal (6--8 in the final thesis).

\textbf{Common mistakes:} - Summarizing papers one by one instead of
synthesizing them thematically - Hypotheses that appear out of nowhere
without a theoretical build-up - Inconsistent naming of constructs (use
the same term throughout) - Too many or too few hypotheses --- each
hypothesis needs a full theoretical argument

\begin{tcolorbox}[enhanced jigsaw, bottomrule=.15mm, opacityback=0, toprule=.15mm, colframe=quarto-callout-tip-color-frame, left=2mm, rightrule=.15mm, title=\textcolor{quarto-callout-tip-color}{\faLightbulb}\hspace{0.5em}{The cone structure}, colbacktitle=quarto-callout-tip-color!10!white, bottomtitle=1mm, colback=white, breakable, toptitle=1mm, titlerule=0mm, leftrule=.75mm, arc=.35mm, opacitybacktitle=0.6, coltitle=black]

Each section of your theory chapter should start broad (definitions,
background) and narrow toward your specific hypotheses. Think of it as a
funnel.

\end{tcolorbox}

\subsection{3. Hypotheses (or
sub-questions)}\label{hypotheses-or-sub-questions-1}

State all hypotheses clearly and separately. Each hypothesis: - Proposes
exactly \textbf{one effect} (one direction, one pair of variables) - Is
numbered (H1, H2, H3 etc.; sub-hypotheses as H1a, H1b) - Uses the
\textbf{same variable names} as your theory, conceptual model, and
method - Is supported by a clear theoretical argument in the theory
chapter

For \textbf{moderation}: \emph{``The effect of X on Y will be moderated
by Z, such that the effect will be stronger/weaker when Z is high.''}

For \textbf{mediation}: \emph{``The effect of X on Y will be mediated by
M.''}

For \textbf{qualitative research}: state sub-questions instead. Each
sub-question should address one distinct aspect of the main RQ.

\subsection{4. Method}\label{method}

The method section describes how you will answer your research question.
At the proposal stage, it should be as detailed as possible, even if
some specifics will be refined later.

\subsubsection{Participants / data
source}\label{participants-data-source}

\begin{itemize}
\tightlist
\item
  Who are the participants (or what is the corpus)?
\item
  How will you recruit them (or sample the data)?
\item
  What is the intended sample size? Include a \textbf{power analysis} or
  justify your \emph{N} (e.g., for experimental research: a medium
  effect size with 80\% power for a between-subjects ANOVA requires
  \textasciitilde52 per cell)
\item
  For qualitative research: how many participants, using what sampling
  logic?
\end{itemize}

\subsubsection{Design}\label{design}

\begin{itemize}
\tightlist
\item
  State the type of design clearly
\item
  For experiments: write out the full design notation (e.g., \emph{``a 2
  {[}X: condition A vs.~condition B{]} × 2 {[}Y: high vs.~low{]}
  between-subjects experiment''})
\item
  Name all independent, dependent, and control variables
\end{itemize}

\subsubsection{Procedure}\label{procedure}

\begin{itemize}
\tightlist
\item
  Describe the study sequence step by step
\item
  For experiments: describe exactly what each condition will look like
\end{itemize}

\subsubsection{Measures /
operationalization}\label{measures-operationalization}

For each variable: - Name the scale or instrument - Cite the original
source - State the number of items and the response format - Note
reliability from the original source (if available)

A first draft of the \textbf{questionnaire, experimental materials, or
codebook} can be included as an appendix to the proposal.

\subsection{5. Analysis plan}\label{analysis-plan}

The analysis plan is one of the most important parts of the proposal and
is often where students struggle. For \textbf{each hypothesis or
sub-question}, specify:

\begin{itemize}
\tightlist
\item
  Which statistical test you will use
\item
  Which variables are included (exactly as named in the method section)
\item
  What the criteria are for accepting or rejecting the hypothesis
\end{itemize}

\textbf{Example analysis plan entries:}

\begin{quote}
\emph{H1 (main effect of framing on policy support) will be tested with
an independent-samples t-test, with news framing condition (political
vs.~scientific) as the independent variable and policy support (scale
average) as the dependent variable. H1 will be supported if the effect
is statistically significant (p \textless{} .05) and in the hypothesized
direction.}
\end{quote}

\begin{quote}
\emph{H3 (moderation of framing effect by political identity) will be
tested with a 2-step hierarchical regression. In Step 1, news framing
(dummy-coded) and political identity are entered as predictors of policy
support. In Step 2, the interaction term (framing × political identity)
is added. H3 will be supported if the interaction term in Step 2 is
statistically significant.}
\end{quote}

\begin{quote}
\emph{H5 (mediation of X on Y via M) will be tested using Hayes' PROCESS
macro (Model 4) with 5,000 bootstrap samples. The indirect effect will
be considered significant if the 95\% bootstrap confidence interval does
not include zero.}
\end{quote}

\textbf{For qualitative research:} describe your analytical approach
(e.g., thematic analysis following Braun \& Clarke, 2006), how you will
code the data, and how you will ensure trustworthiness (e.g., member
checking, peer debriefing).

\begin{tcolorbox}[enhanced jigsaw, bottomrule=.15mm, opacityback=0, toprule=.15mm, colframe=quarto-callout-important-color-frame, left=2mm, rightrule=.15mm, title=\textcolor{quarto-callout-important-color}{\faExclamation}\hspace{0.5em}{The analysis plan is binding}, colbacktitle=quarto-callout-important-color!10!white, bottomtitle=1mm, colback=white, breakable, toptitle=1mm, titlerule=0mm, leftrule=.75mm, arc=.35mm, opacitybacktitle=0.6, coltitle=black]

Your supervisor approves your analysis plan. Deviating from it
substantially in the final thesis without justification is a problem. If
circumstances require a change (e.g., your manipulation fails), discuss
this with your supervisor immediately.

\end{tcolorbox}

\section{Getting to ``sufficient''}\label{getting-to-sufficient}

Your supervisor evaluates the proposal primarily on:

\begin{longtable}[]{@{}
  >{\raggedright\arraybackslash}p{(\columnwidth - 2\tabcolsep) * \real{0.5000}}
  >{\raggedright\arraybackslash}p{(\columnwidth - 2\tabcolsep) * \real{0.5000}}@{}}
\toprule\noalign{}
\begin{minipage}[b]{\linewidth}\raggedright
Criterion
\end{minipage} & \begin{minipage}[b]{\linewidth}\raggedright
What they are looking for
\end{minipage} \\
\midrule\noalign{}
\endhead
\bottomrule\noalign{}
\endlastfoot
\textbf{Feasibility} & Can you realistically do this study within the
timeframe? \\
\textbf{Fit between RQ and design} & Does the design actually answer the
research question? \\
\textbf{Theoretical grounding} & Are the hypotheses derived from solid
theoretical argumentation? \\
\textbf{Analysis plan clarity} & Do you know exactly which test to run
for each hypothesis? \\
\textbf{Consistency} & Are the same variable names used throughout
theory, method, and analysis plan? \\
\end{longtable}

\section{Practical tips}\label{practical-tips-1}

\begin{itemize}
\tightlist
\item
  \textbf{Don't wait for the perfect idea.} Submit the best proposal you
  can. The proposal is a plan, not a guarantee. It can be refined.
\item
  \textbf{Be very specific about your design.} Vague descriptions of ``a
  survey'' or ``an experiment'' are not enough. What exact manipulation?
  What exact measures?
\item
  \textbf{Name your scales.} ``I will measure anxiety'' is not
  sufficient. ``I will measure anxiety using the 6-item STAI-6 (Marteau
  \& Bekker, 1992)'' is.
\item
  \textbf{Map your analysis plan against your hypotheses one by one.} If
  you have 5 hypotheses, you need 5 analysis plan entries --- no more,
  no less.
\item
  \textbf{Process peer feedback before submitting.} Fellow students can
  often spot inconsistencies in variable names or gaps in the
  theoretical argument that you miss when reading your own work.
\end{itemize}

\part{The Research Process}

\chapter{From Topic to Research
Question}\label{from-topic-to-research-question}

\section{Starting with a topic}\label{starting-with-a-topic}

A thesis begins with a \textbf{topic of interest} --- a broad area
within communication science that you want to explore. Your supervisor
offers research themes within their expertise. Working within your
supervisor's theme ensures your project is relevant and that you benefit
fully from their expertise.

If you propose your own topic, make sure it is:

\begin{itemize}
\tightlist
\item
  \textbf{Interesting and feasible} within the timeframe
\item
  \textbf{Anchored in existing literature} (you must be able to find a
  research gap)
\item
  \textbf{Researchable} with data you can realistically collect
\item
  Something a willing supervisor can guide
\end{itemize}

\section{Identifying the research
gap}\label{identifying-the-research-gap}

A good research question emerges from a \textbf{gap in the literature}.
After reading the relevant literature in your area, ask yourself:

\begin{itemize}
\tightlist
\item
  What do we know already?
\item
  What is still unclear, inconsistent, or unexplored?
\item
  Why does this gap matter --- theoretically and/or for society?
\end{itemize}

The gap is your justification for doing the study. It should be clearly
articulated in your Introduction and elaborated in your Theory chapter.

\section{Formulating the research
question}\label{formulating-the-research-question}

A good research question is:

\begin{itemize}
\tightlist
\item
  \textbf{Specific} --- it names the constructs/variables you are
  interested in
\item
  \textbf{Researchable} --- it can be answered with empirical data you
  can collect
\item
  \textbf{Theoretically grounded} --- it follows from a gap in existing
  knowledge
\item
  \textbf{Focused} --- one central question, not multiple questions
  combined
\end{itemize}

\begin{tcolorbox}[enhanced jigsaw, bottomrule=.15mm, opacityback=0, toprule=.15mm, colframe=quarto-callout-tip-color-frame, left=2mm, rightrule=.15mm, title=\textcolor{quarto-callout-tip-color}{\faLightbulb}\hspace{0.5em}{Common mistake: too broad}, colbacktitle=quarto-callout-tip-color!10!white, bottomtitle=1mm, colback=white, breakable, toptitle=1mm, titlerule=0mm, leftrule=.75mm, arc=.35mm, opacitybacktitle=0.6, coltitle=black]

``What is the role of social media in society?'' is too broad. A better
version: ``To what extent does exposure to ideologically congruent
social media content predict political polarization among Dutch
adults?''

\end{tcolorbox}

In \textbf{quantitative research}, the main RQ is typically followed by
\textbf{hypotheses} that predict specific relationships. These are
derived from existing theory.

In \textbf{qualitative research} (e.g., grounded theory, thematic
analysis), the main RQ is followed by \textbf{sub-questions} that guide
your data collection and analysis.

\section{Hypotheses}\label{hypotheses}

A hypothesis proposes a specific, testable relationship between
variables. Good hypotheses:

\begin{itemize}
\tightlist
\item
  Propose \textbf{one effect} (not multiple effects in one statement)
\item
  Follow \textbf{logically} from the theoretical argument
\item
  Are supported by \textbf{3--5 previous studies} pointing in the same
  direction
\item
  Use \textbf{consistent terminology} with the theory section,
  conceptual model, and method section
\end{itemize}

Example: \emph{``H1: Higher exposure to health misinformation on social
media will be positively associated with vaccine hesitancy.''}

\begin{tcolorbox}[enhanced jigsaw, bottomrule=.15mm, opacityback=0, toprule=.15mm, colframe=quarto-callout-important-color-frame, left=2mm, rightrule=.15mm, title=\textcolor{quarto-callout-important-color}{\faExclamation}\hspace{0.5em}{Moderation and mediation in separate hypotheses}, colbacktitle=quarto-callout-important-color!10!white, bottomtitle=1mm, colback=white, breakable, toptitle=1mm, titlerule=0mm, leftrule=.75mm, arc=.35mm, opacitybacktitle=0.6, coltitle=black]

Main effects and moderation/mediation effects must be stated in separate
hypotheses. Do not combine them.

\end{tcolorbox}

\section{Sub-questions (qualitative
research)}\label{sub-questions-qualitative-research}

In qualitative projects, the main RQ is broken down into
\textbf{sub-questions} that:

\begin{itemize}
\tightlist
\item
  Follow logically from the main RQ
\item
  Address specific aspects or steps of your research
\item
  Are grounded in 3--5 previous studies
\item
  Are formulated distinctively (each addresses a clearly different
  focus)
\end{itemize}

Each sub-question ends with a \textbf{single question mark} --- a common
mistake is combining multiple questions before one question mark.

\section{The conceptual model}\label{the-conceptual-model}

For quantitative research, your Theory chapter ends with a
\textbf{conceptual model} --- a visual diagram that maps out the
constructs and their hypothesized relationships. Every variable
mentioned in your hypotheses must appear in the model. Every variable in
the model must appear in your hypotheses and method section.

Use consistent terminology throughout: the same construct should have
the same name in the introduction, theory, conceptual model, method, and
results.

\chapter{Research Design}\label{research-design-1}

Choosing the right research design is one of the most important
decisions you make in your thesis. Your design must fit your research
question. This chapter outlines the main options available to you.

\section{Quantitative vs.~qualitative
research}\label{quantitative-vs.-qualitative-research}

\begin{longtable}[]{@{}
  >{\raggedright\arraybackslash}p{(\columnwidth - 4\tabcolsep) * \real{0.3333}}
  >{\raggedright\arraybackslash}p{(\columnwidth - 4\tabcolsep) * \real{0.3333}}
  >{\raggedright\arraybackslash}p{(\columnwidth - 4\tabcolsep) * \real{0.3333}}@{}}
\toprule\noalign{}
\begin{minipage}[b]{\linewidth}\raggedright
\end{minipage} & \begin{minipage}[b]{\linewidth}\raggedright
Quantitative
\end{minipage} & \begin{minipage}[b]{\linewidth}\raggedright
Qualitative
\end{minipage} \\
\midrule\noalign{}
\endhead
\bottomrule\noalign{}
\endlastfoot
\textbf{Goal} & Test hypotheses about relationships between variables &
Explore meanings, experiences, patterns \\
\textbf{Data} & Numerical (survey responses, coded content) & Text,
images, audio (interviews, documents) \\
\textbf{Analysis} & Statistical tests (regression, ANOVA, etc.) &
Thematic analysis, grounded theory, discourse analysis \\
\textbf{Output} & Generalizable findings across populations & Rich,
contextualized insights \\
\textbf{Thesis structure} & Hypotheses + conceptual model &
Sub-questions \\
\end{longtable}

Both approaches are valid in Communication Science. The choice depends
on your research question and the state of knowledge in your area.

\section{Survey research}\label{survey-research}

A survey measures variables via a questionnaire administered to a sample
of participants. It is the most common design in Communication Science
master's theses.

\textbf{When to use it:} When you want to assess the prevalence of
variables and/or examine associations and relationships between them in
a larger sample.

\textbf{Key considerations:} - Sample size and recruitment (specify
power analysis rationale) - Validated scales vs.~self-constructed items
- Platform (Qualtrics is the standard at VU) - Online convenience
sampling vs.~probability sampling

\section{Experimental research}\label{experimental-research}

An experiment manipulates one or more independent variables and measures
effects on dependent variables, with random assignment of participants
to conditions.

\textbf{When to use it:} When you want to establish causal effects.

\textbf{Design notation:} Describe your design clearly, e.g.:

\begin{quote}
\emph{``A 2 (color of logo: red vs.~blue) × 2 (medium: TV vs.~newspaper)
between-subjects experiment with purchase intention as the dependent
variable.''}
\end{quote}

Always specify: - Whether it is between-subjects, within-subjects, or
mixed - Whether randomization was used (it should be) - Whether you
conducted a \textbf{manipulation check}

\section{Content analysis}\label{content-analysis}

Content analysis systematically codes and analyzes a corpus of media
content (text, images, video, social media posts).

\textbf{Types:} - \textbf{Manual content analysis}: human coders apply a
codebook; requires intercoder reliability testing - \textbf{Automated
content analysis}: computational methods (e.g., dictionary-based,
supervised machine learning, LLMs) via R, AmCat, or Coosto

\textbf{Key considerations:} - How was the corpus assembled? (sampling
strategy) - What is the unit of analysis? - How were variables
operationalized and coded? - What are the reliability and validity of
your coding?

\section{Qualitative research}\label{qualitative-research}

Qualitative designs explore social phenomena in depth without the goal
of statistical generalization.

\textbf{Common approaches in Communication Science:} -
\textbf{Interviews} (semi-structured, in-depth) - \textbf{Focus groups}
- \textbf{Ethnographic observation}

\textbf{Analysis frameworks:} - \textbf{Thematic analysis} ---
identifying recurring patterns across data - \textbf{Grounded theory}
--- building theory inductively from the data - \textbf{Discourse
analysis} --- analyzing how language constructs meaning

In qualitative theses, the method chapter describes the methodology
(e.g., grounded theory), data collection (who, how many, how recruited),
and data analysis procedures.

\section{The analysis plan}\label{the-analysis-plan}

Your \textbf{Research Proposal must include an analysis plan} that
specifies --- for each hypothesis or sub-question --- exactly which
statistical test on which variables will be used to answer it. For
example:

\begin{quote}
\emph{``H1 will be tested with a simple linear regression with social
media exposure as the predictor and vaccine hesitancy as the outcome
variable.''}
\end{quote}

This plan prevents ad-hoc analysis decisions later. In the final thesis,
the analysis plan is not a separate chapter --- it is integrated into
the Results section.

\begin{tcolorbox}[enhanced jigsaw, bottomrule=.15mm, opacityback=0, toprule=.15mm, colframe=quarto-callout-tip-color-frame, left=2mm, rightrule=.15mm, title=\textcolor{quarto-callout-tip-color}{\faLightbulb}\hspace{0.5em}{Pre-registration}, colbacktitle=quarto-callout-tip-color!10!white, bottomtitle=1mm, colback=white, breakable, toptitle=1mm, titlerule=0mm, leftrule=.75mm, arc=.35mm, opacitybacktitle=0.6, coltitle=black]

Consider pre-registering your analysis plan on the Open Science
Framework (OSF) before data collection. See \href{open-science.qmd}{Open
Science} for details.

\end{tcolorbox}

\chapter{Ethics \& Data Management}\label{ethics-data-management}

\section{The ethical self-check}\label{the-ethical-self-check}

Before you start collecting data, it is \textbf{mandatory} to complete
the \textbf{ethical review self-check} via the Research Ethics Review
Committee (RERC) of the Faculty of Social Sciences (FSW).

You can find the self-check and all relevant information here:
\href{https://vu.nl/en/employee/research-support/research-ethics-review-ssc}{vu.nl/en/employee/research-support/research-ethics-review-ssc}

\begin{tcolorbox}[enhanced jigsaw, bottomrule=.15mm, opacityback=0, toprule=.15mm, colframe=quarto-callout-important-color-frame, left=2mm, rightrule=.15mm, title=\textcolor{quarto-callout-important-color}{\faExclamation}\hspace{0.5em}{Always complete the self-check with your supervisor}, colbacktitle=quarto-callout-important-color!10!white, bottomtitle=1mm, colback=white, breakable, toptitle=1mm, titlerule=0mm, leftrule=.75mm, arc=.35mm, opacitybacktitle=0.6, coltitle=black]

Your supervisor holds final responsibility for the accuracy of your
answers. \textbf{Never complete the self-check without consulting your
supervisor first.} Inaccurate answers can invalidate the ethical
assessment and make your data unusable.

\end{tcolorbox}

\subsection{How the self-check works}\label{how-the-self-check-works}

\begin{enumerate}
\def\labelenumi{\arabic{enumi}.}
\tightlist
\item
  Complete the online survey in consultation with your supervisor
\item
  If your research \textbf{meets} the faculty's ethical guidelines → no
  further review needed. You will receive a PDF confirmation by email.
\item
  If your research \textbf{does not meet} the guidelines (e.g., no
  informed consent, stimuli that may distress participants, or
  vulnerable populations) → further steps are needed
\end{enumerate}

\subsection{If your research does not pass the
self-check}\label{if-your-research-does-not-pass-the-self-check}

\begin{enumerate}
\def\labelenumi{\arabic{enumi}.}
\tightlist
\item
  Discuss the issue with your supervisor --- in many cases, a small
  design adjustment solves the problem
\item
  If the issue cannot be resolved, contact the RERC at rerc.ssc@vu.nl
\item
  Note: thesis projects are generally \textbf{not eligible} for full
  RERC review, unless the project is part of a larger, supervised
  research line that will be published. In that case, the supervisor
  must prepare materials for full review.
\end{enumerate}

\subsection{Practical notes}\label{practical-notes}

\begin{itemize}
\tightlist
\item
  Save the self-check PDF output and store it in your surfdrive folder
  and attach it as an appendix to your thesis
\item
  You can change your answers before the final submission of the form
  --- discuss all answers with your supervisor before pressing submit
\item
  In some cases, your supervisor has already applied for ethical review
  for a research line you are joining. Check with your supervisor
  whether you need to complete the self-check separately.
\end{itemize}

\subsection{Qualtrics: disable personal data
collection}\label{qualtrics-disable-personal-data-collection}

By default, Qualtrics stores personal data (IP address and geolocation).
If you want to keep responses anonymous:

\begin{enumerate}
\def\labelenumi{\arabic{enumi}.}
\tightlist
\item
  Go to \emph{Survey Options} in Qualtrics
\item
  Under \emph{Survey Termination}, enable: \textbf{``Anonymize Response
  --- Do NOT record any personal information and remove contact
  association''}
\end{enumerate}

If you do not do this, you must answer ``yes'' to the question about
collecting personal/sensitive data in the self-check, and you must store
the data securely.

\section{Data management}\label{data-management}

Careful and secure handling of research data is required throughout your
thesis.

\subsection{Storing your data
(surfdrive)}\label{storing-your-data-surfdrive}

The accepted secure method for storing and sharing data with your
supervisor is a \textbf{shared SURFdrive folder}, protected by a
password. Your supervisor creates this folder and grants you access.

Workflow: 1. Place the \textbf{complete raw data file} in the shared
folder with a clear name (e.g.,
\texttt{ALL\_RAW\_DATA\_StudyName\_StudentName.sav}) 2. Make a
\textbf{duplicate} of the raw file --- work only in the copy 3.
\textbf{Remove all personal/sensitive data} from the working file (e.g.,
geolocation, IP addresses, names, email addresses) 4. Clean the working
file and create new variables as needed 5. Keep your \textbf{syntax} (R
or SPSS) updated and annotated throughout

\subsection{Syntax hygiene}\label{syntax-hygiene}

Always work with syntax --- never run analyses from menus without
recording them:

\begin{itemize}
\tightlist
\item
  In \textbf{SPSS}: click ``Paste'' for every command run from a menu;
  add notes with \texttt{*\ comment\ text}
\item
  In \textbf{R}: annotate your script with \texttt{\#\ comment\ text}
\end{itemize}

Your syntax should allow anyone (including you, six months from now) to:

\begin{itemize}
\tightlist
\item
  Re-create the final dataset from the raw data
\item
  Reproduce every reported analysis
\end{itemize}

At the start of your method section, report all cases removed from the
dataset and the reasons (e.g., \emph{``9 cases were excluded because
they did not complete the questionnaire''}).

\begin{tcolorbox}[enhanced jigsaw, bottomrule=.15mm, opacityback=0, toprule=.15mm, colframe=quarto-callout-tip-color-frame, left=2mm, rightrule=.15mm, title=\textcolor{quarto-callout-tip-color}{\faLightbulb}\hspace{0.5em}{Syntax is part of your grade}, colbacktitle=quarto-callout-tip-color!10!white, bottomtitle=1mm, colback=white, breakable, toptitle=1mm, titlerule=0mm, leftrule=.75mm, arc=.35mm, opacitybacktitle=0.6, coltitle=black]

Reproducible syntax is a sign of scientific rigor. Your supervisor needs
access to your data and syntax (via surfdrive) when assessing your
thesis.

\end{tcolorbox}

\part{Writing Your Thesis}

\chapter{Structure \& Formalia}\label{structure-formalia}

\section{Length and formatting}\label{length-and-formatting}

\begin{longtable}[]{@{}
  >{\raggedright\arraybackslash}p{(\columnwidth - 2\tabcolsep) * \real{0.5000}}
  >{\raggedright\arraybackslash}p{(\columnwidth - 2\tabcolsep) * \real{0.5000}}@{}}
\toprule\noalign{}
\begin{minipage}[b]{\linewidth}\raggedright
\end{minipage} & \begin{minipage}[b]{\linewidth}\raggedright
Requirement
\end{minipage} \\
\midrule\noalign{}
\endhead
\bottomrule\noalign{}
\endlastfoot
\textbf{Minimum length} & 8,000 words \\
\textbf{Maximum length} & 9,000 words (10\% margin allowed = up to
9,900) \\
\textbf{Word count excludes} & Title page, abstract, references, tables,
appendices, and (for qualitative research) citations from data \\
\textbf{Line spacing} & 1.5 \\
\textbf{Font} & 12pt Times New Roman \\
\textbf{Pages (approx.)} & 22--28 \\
\textbf{Word count location} & State on title page \\
\end{longtable}

\begin{tcolorbox}[enhanced jigsaw, bottomrule=.15mm, opacityback=0, toprule=.15mm, colframe=quarto-callout-note-color-frame, left=2mm, rightrule=.15mm, title=\textcolor{quarto-callout-note-color}{\faInfo}\hspace{0.5em}{Length is about quality, not padding}, colbacktitle=quarto-callout-note-color!10!white, bottomtitle=1mm, colback=white, breakable, toptitle=1mm, titlerule=0mm, leftrule=.75mm, arc=.35mm, opacitybacktitle=0.6, coltitle=black]

If your project is genuinely complex and cannot be reported within 9,000
words, discuss this with your supervisor. However, being overly long
without justification may negatively affect your grade. Less-important
material can usually be moved to an appendix.

\end{tcolorbox}

\section{Required structure}\label{required-structure}

Your thesis must follow this structure. The page counts below are
guidelines, not strict requirements.

\begin{longtable}[]{@{}
  >{\raggedright\arraybackslash}p{(\columnwidth - 4\tabcolsep) * \real{0.1304}}
  >{\raggedright\arraybackslash}p{(\columnwidth - 4\tabcolsep) * \real{0.3913}}
  >{\raggedright\arraybackslash}p{(\columnwidth - 4\tabcolsep) * \real{0.4783}}@{}}
\toprule\noalign{}
\begin{minipage}[b]{\linewidth}\raggedright
\#
\end{minipage} & \begin{minipage}[b]{\linewidth}\raggedright
Section
\end{minipage} & \begin{minipage}[b]{\linewidth}\raggedright
Guideline
\end{minipage} \\
\midrule\noalign{}
\endhead
\bottomrule\noalign{}
\endlastfoot
0 & \textbf{Title page} & --- \\
1 & \textbf{Table of contents} & --- \\
2 & \textbf{Abstract / Summary} & \textasciitilde0.5 page \\
3 & \textbf{Introduction} & 1--2 pages \\
4 & \textbf{Theory} & 6--8 pages \\
5 & \textbf{Method} & 2--3 pages \\
6 & \textbf{Results} & 3--4 pages (qualitative: 6--8 pages) \\
7 & \textbf{Conclusion and Discussion} & 4--5 pages \\
8 & \textbf{Reference list} & ≥ 30 references \\
9 & \textbf{Statement about Generative AI and AI-Assisted Technology
usage} & After references; not in word count \\
10 & \textbf{Appendices} & --- \\
\end{longtable}

Depending on your research design, you may deviate from this structure
--- but always consult your supervisor.

\subsection{Title page}\label{title-page}

The title page should include at minimum: - Title of your thesis - Your
name and student number - Course name and code (S\_MTcw) - Supervisor's
name - Date of submission - Word count of the main text - VU Amsterdam /
Faculty of Social Sciences / Department of Communication Science

\subsection{Abstract}\label{abstract}

The abstract (approximately half a page) is written \textbf{last} but
appears first. It concisely describes: - The reason for the study and
the research question/hypotheses - The method - The key results - The
main conclusion

\subsection{Appendices}\label{appendices}

At minimum, include: - Your questionnaire and/or stimulus materials -
Your codebook (for content analysis) - The RERC ethical self-check PDF

Variable names in the appendices should match exactly those used in your
Method and Results sections.

\section{Formatting (APA)}\label{formatting-apa}

All text, tables, figures, the title page, table of contents, abstract,
and reference list must follow \textbf{APA 7th edition} guidelines. See
\href{citation.qmd}{Citing \& Referencing} for details and
\href{https://apastyle.apa.org/}{apastyle.apa.org} for the full style
guide.

Key formatting rules: - Statistical symbols are italicized: \emph{M},
\emph{SD}, \emph{F}, \emph{t}, \emph{p}, \emph{r}, \emph{N}, \emph{n} -
All numbers rounded to \textbf{two decimal places} (except \emph{p}
\textless{} .001) - No leading zero on \emph{p} values and correlations:
write \emph{p} = .03, not \emph{p} = 0.03 - Report exact \emph{p}
values: \emph{p} = .03, except when \emph{p} \textless{} .001 - Spaces
after = and \textless{} signs - Tables and figures must include all
necessary information and labels

\section{AI disclosure statement}\label{ai-disclosure-statement}

All theses must include a \textbf{Statement about Generative AI and
AI-Assisted Technology usage} after the reference list (not counted in
the word count). See \href{ai-tools.qmd}{Using AI Tools} for what this
statement must contain.

\section{Submission}\label{submission}

Submit your final thesis: 1. \textbf{On Canvas} (where it will be
checked for plagiarism) 2. \textbf{By email} to both your supervisor and
second reviewer (check with your supervisor whether a paper version is
also required) 3. \textbf{Ensure your supervisor has access} to your
final data file and syntax via the shared surfdrive folder

After receiving a passing grade, you must also \textbf{upload your
thesis to the UBVU thesis database} --- this is a prerequisite for
receiving your diploma. See \href{assessment.qmd}{Assessment \&
Graduation} for details.

\chapter{Introduction \& Theory
Chapter}\label{introduction-theory-chapter}

\section{The Abstract}\label{the-abstract}

Although it appears first in your thesis, the abstract is the
\textbf{last section you write}. It concisely summarizes:

\begin{itemize}
\tightlist
\item
  The reason for the study and the research question or hypotheses
\item
  The method used
\item
  The key findings
\item
  The main conclusion
\end{itemize}

Keep it to approximately half a page. Focus on the most important points
--- do not include citations.

\section{The Introduction}\label{the-introduction}

The introduction presents your topic and frames the study. A good
introduction is \textbf{concise} and focuses on two questions:
\emph{``What has been done in this area?''} and \emph{``Why is the
current study useful?''}

\subsection{Standard structure}\label{standard-structure}

\begin{enumerate}
\def\labelenumi{\arabic{enumi}.}
\tightlist
\item
  \textbf{Reason} (\emph{aanleiding}) --- Why is this topic relevant
  now? What social or scientific development motivates this study?
\item
  \textbf{Topic introduction and knowledge gap} --- Briefly introduce
  the topic and name the gap in existing research (keep it short here;
  details follow in the Theory chapter)
\item
  \textbf{Problem statement} --- Articulate the specific research
  problem
\item
  \textbf{Research question} --- State the central RQ explicitly and
  clearly
\item
  \textbf{Societal and scientific relevance} --- Explain why answering
  this RQ matters
\end{enumerate}

For \textbf{qualitative research}, also briefly summarize your
methodology (e.g., grounded theory) and methods (e.g., interviews,
thematic analysis) at the end of the introduction.

\subsection{Checklist for the
introduction}\label{checklist-for-the-introduction}

\begin{itemize}
\tightlist
\item[$\square$]
  Is a clear motivation for the research problem provided?
\item[$\square$]
  Is the research problem clearly delineated?
\item[$\square$]
  Is the research gap clearly identified?
\item[$\square$]
  Are central concepts clearly defined and consistently named throughout
  the thesis?
\item[$\square$]
  Is the research question stated explicitly and clearly?
\item[$\square$]
  Is the research question researchable?
\item[$\square$]
  Is both the \textbf{societal} and \textbf{scientific} relevance
  demonstrated?
\end{itemize}

\section{The Theory Chapter}\label{the-theory-chapter}

The theory chapter reviews the relevant literature and major theories,
shows how your research builds on existing work, and leads to your
hypotheses (or sub-questions for qualitative research).

\subsection{Structure}\label{structure}

A theory chapter typically:

\begin{enumerate}
\def\labelenumi{\arabic{enumi}.}
\tightlist
\item
  Defines and explains the central constructs/concepts
\item
  Builds a theoretical argument for the relationships between constructs
\item
  Discusses relevant prior empirical findings
\item
  Culminates in clearly formulated hypotheses (or sub-questions) and a
  conceptual model
\end{enumerate}

Each sub-section should follow a \textbf{cone shape}: start with general
definitions and background, then narrow to specific relationships
relevant to your hypotheses.

\subsection{Content checklist}\label{content-checklist}

\begin{itemize}
\tightlist
\item[$\square$]
  Is the theory relevant to the research question?
\item[$\square$]
  Is the literature up-to-date (primarily recent peer-reviewed journal
  articles)?
\item[$\square$]
  Is a critical stance taken (noting relationships, contradictions,
  gaps)?
\item[$\square$]
  Are highly relevant studies discussed in appropriate detail (including
  their methods and findings)?
\item[$\square$]
  Is the labeling of variables/constructs consistent throughout?
\end{itemize}

\subsection{Structure and writing
checklist}\label{structure-and-writing-checklist}

\begin{itemize}
\tightlist
\item[$\square$]
  Is there a clear logical thread (\emph{rode lijn}) throughout?
\item[$\square$]
  Does each section begin with a topical sentence?
\item[$\square$]
  Are sections balanced in length?
\item[$\square$]
  Is the conceptual model (for quantitative research) included at the
  end?
\item[$\square$]
  Is every construct in the hypotheses represented in the conceptual
  model?
\end{itemize}

\subsection{Vocabulary and references}\label{vocabulary-and-references}

\begin{itemize}
\tightlist
\item[$\square$]
  Is the vocabulary academic (avoiding contractions and colloquialisms)?
\item[$\square$]
  Is every claim supported by a reference?
\item[$\square$]
  Are in-text citations in correct APA format?
\item[$\square$]
  Are direct quotations justified, and do they include a page number?
\end{itemize}

\subsection{Hypotheses checklist}\label{hypotheses-checklist}

\begin{itemize}
\tightlist
\item[$\square$]
  Does each hypothesis propose one effect only (not multiple)?
\item[$\square$]
  Are moderation and mediation effects in separate hypotheses?
\item[$\square$]
  Are terms used consistently across theory, conceptual model, and
  method?
\item[$\square$]
  Is each hypothesis supported by 3--5 prior studies?
\end{itemize}

\begin{tcolorbox}[enhanced jigsaw, bottomrule=.15mm, opacityback=0, toprule=.15mm, colframe=quarto-callout-tip-color-frame, left=2mm, rightrule=.15mm, title=\textcolor{quarto-callout-tip-color}{\faLightbulb}\hspace{0.5em}{The golden thread}, colbacktitle=quarto-callout-tip-color!10!white, bottomtitle=1mm, colback=white, breakable, toptitle=1mm, titlerule=0mm, leftrule=.75mm, arc=.35mm, opacitybacktitle=0.6, coltitle=black]

Your research question should be visible in every paragraph of your
thesis. Whenever you are unsure whether a paragraph belongs, ask
yourself: \emph{does this contribute to answering my research question?}

\end{tcolorbox}

\chapter{Method Chapter}\label{method-chapter}

The method chapter is written in \textbf{past tense} and must be
detailed enough that another researcher could fully \textbf{replicate
your study}. It is highly standardized in structure.

\section{Overall checklist}\label{overall-checklist}

\begin{itemize}
\tightlist
\item[$\square$]
  Is the method section written in past tense?
\item[$\square$]
  Is the level of detail sufficient for replication?
\item[$\square$]
  Is there a clear fit between the method and the
  hypotheses/sub-questions?
\end{itemize}

\section{Quantitative research (survey and
experiment)}\label{quantitative-research-survey-and-experiment}

\subsection{Participants}\label{participants}

Report: - Method of \textbf{recruitment} - \textbf{Demographics}:
frequency and percentage for gender, age, education, etc. -
\textbf{Sample size}, including justification (power analysis) -
\textbf{Exclusions}: exactly how many participants were excluded and for
what reason(s) - The \textbf{final \emph{N}} (must match the \emph{N}
reported in the results section)

\subsection{Design}\label{design-1}

\begin{itemize}
\tightlist
\item
  State the \textbf{type of design} clearly: survey, experiment,
  between-subjects, within-subjects, mixed
\item
  Briefly describe the \textbf{task} (e.g., \emph{``Participants were
  exposed to a news article and then completed a questionnaire''})
\item
  Name the \textbf{independent and dependent variables}
\item
  For experiments: describe factors and conditions clearly using design
  notation (e.g., \emph{``a 2 {[}logo color: red, blue{]} × 2 {[}medium:
  TV, newspaper{]} between-subjects design''}), and confirm
  randomization
\end{itemize}

\subsection{Procedure}\label{procedure-1}

Describe in \textbf{chronological order}: - What participants were
exposed to (especially condition differences in experiments) - What
tasks and measurements were completed - Key instructions given (purpose
of study, debriefing)

\subsection{Materials}\label{materials}

\begin{itemize}
\tightlist
\item
  Describe \textbf{experimental stimuli} or materials, including exact
  differences between conditions
\item
  All materials must be provided in the appendix
\end{itemize}

\subsection{Measures}\label{measures}

For each variable, report: - Variable name (in \textbf{bold} or
\emph{italics}) - What it measures - Source reference (original scale) -
Number of items - Example item(s) --- or all items if few; refer to
appendix for longer scales - Response format (e.g., \emph{``5-point
Likert scale from 1 = strongly disagree to 5 = strongly agree''}) -
\emph{M} and \emph{SD} in your sample, along with minimum and maximum
values - Reliability (Cronbach's α) - How the final score was computed
(e.g., average of items; recoded items) - Direction (e.g.,
\emph{``Higher scores indicate greater vaccine hesitancy''}) - For
experiments: a \textbf{manipulation check}

Organize measures under subheadings: Independent variables, Dependent
variables, Control variables.

\begin{tcolorbox}[enhanced jigsaw, bottomrule=.15mm, opacityback=0, toprule=.15mm, colframe=quarto-callout-warning-color-frame, left=2mm, rightrule=.15mm, title=\textcolor{quarto-callout-warning-color}{\faExclamationTriangle}\hspace{0.5em}{Consistency is critical}, colbacktitle=quarto-callout-warning-color!10!white, bottomtitle=1mm, colback=white, breakable, toptitle=1mm, titlerule=0mm, leftrule=.75mm, arc=.35mm, opacitybacktitle=0.6, coltitle=black]

The variable names in your method section must be \textbf{identical} to
those used in your hypotheses, conceptual model, and results section.
This is one of the most common sources of errors in MA theses.

\end{tcolorbox}

\section{Content analysis}\label{content-analysis-1}

The method chapter for a content analysis includes:

\begin{itemize}
\tightlist
\item
  A short description of the \textbf{research design}
\item
  \textbf{Data collection}: how the corpus was assembled (sampling
  strategy, date range, sources)
\item
  \textbf{Content analysis method}: manual (intercoder reliability) or
  automated (tool used, e.g., R, AmCat, Coosto)
\item
  \textbf{Operationalization}: search terms or codebook
\item
  \textbf{Reliability}: intercoder agreement or precision/recall of
  search terms
\end{itemize}

\section{Qualitative research}\label{qualitative-research-1}

The method chapter for qualitative research describes:

\begin{itemize}
\tightlist
\item
  \textbf{Methodology}: the chosen framework (e.g., grounded theory) and
  its rationale
\item
  \textbf{Case study} (if applicable): why this case is appropriate for
  the research question
\item
  \textbf{Participants/key informants}: recruitment method,
  demographics, sample size, inclusion/exclusion criteria
\item
  \textbf{Data collection}: method (interviews, observation, etc.),
  materials used (topic guide, prompts), where conducted
\item
  \textbf{Data analysis}: analytical approach (thematic analysis,
  grounded theory coding, discourse analysis), analytical materials
\end{itemize}

All choices should be supported by methodological literature with proper
APA citations.

\chapter{Results Chapter}\label{results-chapter}

The results section is a \textbf{matter-of-fact description} of your
findings. It is written in \textbf{past tense} and focuses on answering
your research questions and testing your hypotheses. Extensive
interpretation belongs in the Discussion --- here you report what you
found.

That said, good results sections are well-structured and tell a coherent
``story'' that guides the reader through the analyses step by step.

\section{Structure for quantitative
research}\label{structure-for-quantitative-research}

Report sections in this order:

\subsection{1. Descriptive statistics}\label{descriptive-statistics}

Report means (\emph{M}), standard deviations (\emph{SD}), and
correlations among measured variables. Present these in a table or
figure. In the text, describe the main and noteworthy correlations. Note
whether control variables correlate with (in)dependent variables --- if
so, you may want to include them in follow-up analyses.

\subsection{2. Manipulation check (for
experiments)}\label{manipulation-check-for-experiments}

Verify that the experimental manipulation worked as intended. Also check
randomization: are demographic variables (gender, age, education)
balanced across conditions? Report these checks even if no significant
differences were found.

\subsection{3. Hypothesis testing}\label{hypothesis-testing}

For each hypothesis: 1. \textbf{Restate} the hypothesis briefly 2.
\textbf{Name the test} used and which variables are included (e.g.,
\emph{``We conducted a linear regression with X as predictor and Y as
outcome''}) 3. \textbf{Report results} with all relevant statistics 4.
\textbf{State a brief conclusion} (e.g., \emph{``H1 was supported''} or
\emph{``H1 was not supported''})

Follow the same order as in the Theory chapter. Also briefly report
non-significant findings --- do not omit them.

\subsection{4. Secondary / exploratory analyses
(optional)}\label{secondary-exploratory-analyses-optional}

If you want to explore the data further (e.g., covariate analyses with
control variables), include these in a clearly labeled separate section.

\section{APA statistical reporting}\label{apa-statistical-reporting}

Follow APA 7th edition formatting:

\begin{itemize}
\tightlist
\item
  Statistical symbols are italicized: \emph{M}, \emph{SD}, \emph{N},
  \emph{n}, \emph{F}, \emph{t}, \emph{r}, \emph{p}, \emph{b}, \emph{β},
  \emph{χ²}
\item
  Round all numbers to \textbf{two decimal places}
\item
  No leading zero on \emph{p} values and correlations: \emph{p} = .03,
  \emph{r} = .45 (not 0.03, 0.45)
\item
  Report \textbf{exact \emph{p} values}: \emph{p} = .03 --- except use
  \emph{p} \textless{} .001 for very small values
\item
  Put spaces after = and \textless{}
\item
  Report \textbf{effect sizes} (Cohen's \emph{d}, η², \emph{f²}, etc.)
\end{itemize}

\textbf{Examples:}

\begin{quote}
The regression analysis revealed a significant positive effect of social
media exposure on vaccine hesitancy, \emph{b} = 0.34, \emph{SE} = 0.08,
\emph{t}(298) = 4.25, \emph{p} \textless{} .001, \emph{β} = .24.
\end{quote}

\begin{quote}
An independent-samples \emph{t}-test showed no significant difference
between conditions, \emph{t}(148) = 1.12, \emph{p} = .264, \emph{d} =
0.18.
\end{quote}

\begin{quote}
The ANOVA yielded a significant main effect of logo color, \emph{F}(1,
196) = 8.34, \emph{p} = .004, η² = .04, but no significant main effect
of medium (\emph{F} \textless{} 1, \emph{ns}).
\end{quote}

\begin{tcolorbox}[enhanced jigsaw, bottomrule=.15mm, opacityback=0, toprule=.15mm, colframe=quarto-callout-important-color-frame, left=2mm, rightrule=.15mm, title=\textcolor{quarto-callout-important-color}{\faExclamation}\hspace{0.5em}{Report all effects, including non-significant ones}, colbacktitle=quarto-callout-important-color!10!white, bottomtitle=1mm, colback=white, breakable, toptitle=1mm, titlerule=0mm, leftrule=.75mm, arc=.35mm, opacitybacktitle=0.6, coltitle=black]

Non-significant findings are still findings. Report them briefly:
\emph{``We observed no main effect of {[}variable{]} (F \textless{} 1,
ns)''}.

\end{tcolorbox}

\section{Tables and figures}\label{tables-and-figures}

\begin{itemize}
\tightlist
\item
  Use tables and figures only when they \textbf{add information} that
  cannot be conveyed efficiently in text
\item
  Format all tables in \textbf{APA style} (not copy-pasted from SPSS
  output)
\item
  Every table and figure needs:

  \begin{itemize}
  \tightlist
  \item
    A title (above for tables, below for figures)
  \item
    Clear axis labels or column headers
  \item
    A note if abbreviations are used
  \end{itemize}
\item
  Variable labels in tables must match those used in the method section
\end{itemize}

\section{Results in qualitative
research}\label{results-in-qualitative-research}

In qualitative results, the chapter demonstrates your analysis and
presents empirical findings. Structure it around your
\textbf{sub-questions} (usually resulting in three main sections). For
each sub-question:

\begin{enumerate}
\def\labelenumi{\arabic{enumi}.}
\tightlist
\item
  Restate the sub-question
\item
  Present the analytical findings (from selective coding)
\item
  Support findings with representative quotes from the data
\item
  Briefly conclude
\end{enumerate}

\textbf{Quotes} serve as illustrations --- not as substitutes for
analysis. Contextualize each quote. Format quotes in italics, indented.
Use a respondent ID rather than identifying information.

Avoid discussing theoretical implications extensively in the results
section --- save that for the Discussion.

\chapter{Conclusion \& Discussion}\label{conclusion-discussion}

The conclusion and discussion chapter is the final major section of your
thesis. It is where you step back from the data and reflect on the
meaning, implications, and limitations of your findings.

\begin{tcolorbox}[enhanced jigsaw, bottomrule=.15mm, opacityback=0, toprule=.15mm, colframe=quarto-callout-tip-color-frame, left=2mm, rightrule=.15mm, title=\textcolor{quarto-callout-tip-color}{\faLightbulb}\hspace{0.5em}{Language tense}, colbacktitle=quarto-callout-tip-color!10!white, bottomtitle=1mm, colback=white, breakable, toptitle=1mm, titlerule=0mm, leftrule=.75mm, arc=.35mm, opacitybacktitle=0.6, coltitle=black]

Descriptions of \emph{your} findings → \textbf{past tense} (\emph{``The
results showed that\ldots{}''}) Theoretical implications and conclusions
→ \textbf{present tense} (\emph{``This suggests that\ldots{}''})

\end{tcolorbox}

\section{Standard structure (use max. 3
subheadings)}\label{standard-structure-use-max.-3-subheadings}

\subsection{1. Summary and conclusion}\label{summary-and-conclusion}

\textbf{Opening paragraph:} - Restate the main focus and purpose of the
study - Briefly mention the approach (design) - Give a clear, concise
overall summary of the answer to your research question

\textbf{Middle part --- detailed summary and discussion of findings:}

For each hypothesis or sub-question: - Summarize the finding - Relate it
to the \textbf{specific literature} used to motivate the hypothesis -
Explain whether it is consistent or inconsistent with prior work --- and
\emph{why} - For unexpected findings: introduce possible explanations,
potentially drawing on additional literature - Offer multiple
interpretations where relevant - Make cross-connections between findings

Common mistakes to avoid: - Saying \emph{``This finding is in line with
previous research (ref1, ref2)''} --- explain \emph{how} and \emph{why}
- Ignoring non-significant findings --- discuss them and consider
whether the study was adequately powered

\textbf{Middle part --- broader implications:} - \textbf{Theoretical
implications}: What does your study contribute to theory? Does it
extend, refine, or challenge existing frameworks? - \textbf{Practical
implications}: What do your findings mean for practitioners,
policymakers, or society? - Connect back to the relevance you
established in the Introduction

\textbf{Closing of the first section:} - Consider the broader picture
--- what has your study contributed to the field? - Do the findings
yield new insights?

\subsection{2. Strengths, limitations, and future
research}\label{strengths-limitations-and-future-research}

\textbf{Strengths:} - What methodological or conceptual strengths does
your study have? (These give meaning to the limitations --- there are
always trade-offs)

\textbf{Limitations:} - Be thorough and honest - Address methodological
limitations (e.g., sampling, measurement, generalizability) - For
qualitative research: discuss applicability of findings (situations,
variations)

\textbf{Future research:} - Suggestions must go beyond just \emph{fixing
the limitations} - Propose new, theoretically interesting research
questions that follow from your conclusions - These suggestions should
be grounded in your theoretical implications

\subsection{3. Closing paragraph / final
conclusion}\label{closing-paragraph-final-conclusion}

End with a strong, clear paragraph: - Concisely restate the goal of the
study - State a clear and well-founded conclusion - Provide a direct
answer to the research question and hypotheses - End with a strong
take-home message

\begin{tcolorbox}[enhanced jigsaw, bottomrule=.15mm, opacityback=0, toprule=.15mm, colframe=quarto-callout-warning-color-frame, left=2mm, rightrule=.15mm, title=\textcolor{quarto-callout-warning-color}{\faExclamationTriangle}\hspace{0.5em}{The discussion is the last thing readers read}, colbacktitle=quarto-callout-warning-color!10!white, bottomtitle=1mm, colback=white, breakable, toptitle=1mm, titlerule=0mm, leftrule=.75mm, arc=.35mm, opacitybacktitle=0.6, coltitle=black]

A weak final paragraph is memorable for the wrong reasons. End
confidently with a clear statement of what your study has contributed.

\end{tcolorbox}

\section{Full discussion checklist}\label{full-discussion-checklist}

\begin{itemize}
\tightlist
\item[$\square$]
  Does the opening clearly restate the research focus and approach?
\item[$\square$]
  Does the summary use consistent terminology from the Introduction?
\item[$\square$]
  Is each finding linked to the specific literature that motivated it?
\item[$\square$]
  Are explanations given for both consistent \emph{and} inconsistent
  findings?
\item[$\square$]
  Are cross-connections made between findings?
\item[$\square$]
  Are both theoretical and practical implications discussed?
\item[$\square$]
  Is the limitation section thorough (not just a brief list)?
\item[$\square$]
  Are future research suggestions theoretically grounded?
\item[$\square$]
  Is there a strong closing sentence?
\end{itemize}

\part{Academic Practices}

\chapter{Literature Search}\label{literature-search}

A thorough and systematic literature review is foundational to your
thesis. This chapter offers guidance on how to find, organize, and
manage the literature for your research.

\section{Where to search}\label{where-to-search}

VU Amsterdam provides access to major academic databases. Log in with
your VU credentials at \href{https://ub.vu.nl}{ub.vu.nl}.

\begin{longtable}[]{@{}
  >{\raggedright\arraybackslash}p{(\columnwidth - 2\tabcolsep) * \real{0.5000}}
  >{\raggedright\arraybackslash}p{(\columnwidth - 2\tabcolsep) * \real{0.5000}}@{}}
\toprule\noalign{}
\begin{minipage}[b]{\linewidth}\raggedright
Database
\end{minipage} & \begin{minipage}[b]{\linewidth}\raggedright
Best for
\end{minipage} \\
\midrule\noalign{}
\endhead
\bottomrule\noalign{}
\endlastfoot
\textbf{Scopus} & Broad coverage of social sciences, easy citation
tracking \\
\textbf{Web of Science} & High-quality journals, citation analysis \\
\textbf{PsycINFO} & Psychology and communication research \\
\textbf{Communication Abstracts} & Communication-specific literature \\
\textbf{Google Scholar} & Quick searches, finding grey literature \\
\textbf{JSTOR} & Archive of humanities and social science journals \\
\end{longtable}

For \textbf{qualitative or content analysis} work, additional platforms
like \textbf{LexisNexis} (news content) or \textbf{Coosto/ANP} (social
media) may be relevant.

\section{Search strategies}\label{search-strategies}

\subsection{Keywords and Boolean
operators}\label{keywords-and-boolean-operators}

Construct your search systematically: - Use \textbf{AND} to narrow
results (e.g., \texttt{social\ media\ AND\ misinformation\ AND\ health})
- Use \textbf{OR} to broaden results for synonyms (e.g.,
\texttt{vaccine\ hesitancy\ OR\ vaccine\ refusal}) - Use \textbf{NOT} to
exclude irrelevant terms - Use \textbf{quotation marks} for exact
phrases: \texttt{"filter\ bubble"} - Use \textbf{truncation} with
\texttt{*} for word variations: \texttt{persuad*} finds persuade,
persuasion, persuaded

Keep a record of your search terms, databases used, and results counts
--- you may need to report this in your method section.

\subsection{Snowballing}\label{snowballing}

Once you find a central paper: - \textbf{Forward snowballing}: who has
cited this paper? (use Google Scholar or Scopus) - \textbf{Backward
snowballing}: what does this paper cite?

\subsection{Scoping and narrowing}\label{scoping-and-narrowing}

\begin{itemize}
\tightlist
\item
  Prioritize \textbf{peer-reviewed journal articles} from reputable
  journals
\item
  Prefer \textbf{recent} literature (typically last 10--15 years, unless
  citing foundational work)
\item
  Be selective --- your Theory chapter should discuss only what directly
  contributes to your argument
\end{itemize}

\begin{tcolorbox}[enhanced jigsaw, bottomrule=.15mm, opacityback=0, toprule=.15mm, colframe=quarto-callout-warning-color-frame, left=2mm, rightrule=.15mm, title=\textcolor{quarto-callout-warning-color}{\faExclamationTriangle}\hspace{0.5em}{Quality over quantity}, colbacktitle=quarto-callout-warning-color!10!white, bottomtitle=1mm, colback=white, breakable, toptitle=1mm, titlerule=0mm, leftrule=.75mm, arc=.35mm, opacitybacktitle=0.6, coltitle=black]

The vast majority of your references should be peer-reviewed academic
journal articles. Avoid blogs, Wikipedia, previous BA/MA theses, and PhD
dissertations unless unavoidable.

\end{tcolorbox}

\section{Reference management:
Zotero}\label{reference-management-zotero}

We strongly recommend using
\textbf{\href{https://www.zotero.org}{Zotero}} (free) to manage your
references:

\begin{itemize}
\tightlist
\item
  Install the browser extension to save references directly from
  databases
\item
  Organize references into collections by topic
\item
  Add notes to each reference summarizing relevant content
\item
  Generate in-text citations and reference lists in APA format
\item
  Export \texttt{.bib} files for Quarto, LaTeX, or Word integration
\end{itemize}

\begin{tcolorbox}[enhanced jigsaw, bottomrule=.15mm, opacityback=0, toprule=.15mm, colframe=quarto-callout-tip-color-frame, left=2mm, rightrule=.15mm, title=\textcolor{quarto-callout-tip-color}{\faLightbulb}\hspace{0.5em}{Write as you read}, colbacktitle=quarto-callout-tip-color!10!white, bottomtitle=1mm, colback=white, breakable, toptitle=1mm, titlerule=0mm, leftrule=.75mm, arc=.35mm, opacitybacktitle=0.6, coltitle=black]

When reading a paper, immediately write down the relevant arguments,
findings, and how they connect to your own study. Do not rely on being
able to recall this later. This also prevents you from losing track of
sources.

\end{tcolorbox}

\section{AI tools for literature
search}\label{ai-tools-for-literature-search}

AI tools can support (but not replace) your literature search:

\begin{longtable}[]{@{}
  >{\raggedright\arraybackslash}p{(\columnwidth - 2\tabcolsep) * \real{0.5455}}
  >{\raggedright\arraybackslash}p{(\columnwidth - 2\tabcolsep) * \real{0.4545}}@{}}
\toprule\noalign{}
\begin{minipage}[b]{\linewidth}\raggedright
Tool
\end{minipage} & \begin{minipage}[b]{\linewidth}\raggedright
Use
\end{minipage} \\
\midrule\noalign{}
\endhead
\bottomrule\noalign{}
\endlastfoot
\textbf{Elicit} & AI-powered literature search and summarization \\
\textbf{Research Rabbit} & Discover related papers via citation
networks \\
\textbf{Undermind} & AI-assisted deep literature search \\
\textbf{NotebookLM} (VU-licensed) & Summarize and explore uploaded
documents \\
\textbf{Scite} & Find papers that support or contradict a claim \\
\end{longtable}

\begin{tcolorbox}[enhanced jigsaw, bottomrule=.15mm, opacityback=0, toprule=.15mm, colframe=quarto-callout-important-color-frame, left=2mm, rightrule=.15mm, title=\textcolor{quarto-callout-important-color}{\faExclamation}\hspace{0.5em}{Always verify AI-generated literature}, colbacktitle=quarto-callout-important-color!10!white, bottomtitle=1mm, colback=white, breakable, toptitle=1mm, titlerule=0mm, leftrule=.75mm, arc=.35mm, opacitybacktitle=0.6, coltitle=black]

AI tools can hallucinate references that do not exist or misattribute
findings. \textbf{Always read the original sources yourself} and verify
that citations are real and accurately represented before including them
in your thesis.

\end{tcolorbox}

See \href{ai-tools.qmd}{Using AI Tools} for the full departmental AI
policy.

\chapter{Citing \& Referencing (APA
7th)}\label{citing-referencing-apa-7th}

Your thesis must follow \textbf{APA 7th edition} throughout --- for
in-text citations, the reference list, tables, figures, and statistical
reporting. The authoritative source is
\href{https://apastyle.apa.org/}{apastyle.apa.org}.

\section{In-text citations}\label{in-text-citations}

\subsection{Basic format}\label{basic-format}

\begin{longtable}[]{@{}
  >{\raggedright\arraybackslash}p{(\columnwidth - 4\tabcolsep) * \real{0.3929}}
  >{\raggedright\arraybackslash}p{(\columnwidth - 4\tabcolsep) * \real{0.2857}}
  >{\raggedright\arraybackslash}p{(\columnwidth - 4\tabcolsep) * \real{0.3214}}@{}}
\toprule\noalign{}
\begin{minipage}[b]{\linewidth}\raggedright
Situation
\end{minipage} & \begin{minipage}[b]{\linewidth}\raggedright
Format
\end{minipage} & \begin{minipage}[b]{\linewidth}\raggedright
Example
\end{minipage} \\
\midrule\noalign{}
\endhead
\bottomrule\noalign{}
\endlastfoot
One author & (Author, Year) & (Masur, 2023) \\
Two authors & (Author \& Author, Year) & (Masur \& Teutsch, 2023) \\
Three or more authors & (First author et al., Year) & (Masur et al.,
2023) \\
Narrative citation & Author (Year) found that\ldots{} & Masur (2023)
found that\ldots{} \\
Multiple citations & Alphabetical order, separated by ; & (Jones, 2021;
Masur, 2023) \\
\end{longtable}

\subsection{Direct quotations}\label{direct-quotations}

If you quote directly, you must include a \textbf{page number}:

\begin{quote}
(Masur, 2023, p.~45)
\end{quote}

\begin{tcolorbox}[enhanced jigsaw, bottomrule=.15mm, opacityback=0, toprule=.15mm, colframe=quarto-callout-note-color-frame, left=2mm, rightrule=.15mm, title=\textcolor{quarto-callout-note-color}{\faInfo}\hspace{0.5em}{Minimize direct quotations}, colbacktitle=quarto-callout-note-color!10!white, bottomtitle=1mm, colback=white, breakable, toptitle=1mm, titlerule=0mm, leftrule=.75mm, arc=.35mm, opacitybacktitle=0.6, coltitle=black]

In scientific writing, you should paraphrase and synthesize --- direct
quotations are the exception, not the rule. Only quote when the exact
wording is essential.

\end{tcolorbox}

\subsection{Paraphrasing}\label{paraphrasing}

When you restate an idea in your own words, you still need a citation.
Every factual claim, theoretical argument, or empirical finding that is
not your own must be cited.

\section{Reference list}\label{reference-list}

The reference list appears at the end of the thesis, after the
Conclusion and before the AI statement. Entries are in
\textbf{alphabetical order by first author's last name} and use a
\textbf{hanging indent}.

\subsection{Common reference types}\label{common-reference-types}

\textbf{Journal article:} \textgreater{} Masur, P. K. (2020). How online
privacy literacy affects users' information provision behavior.
\emph{Behaviour \& Information Technology}, \emph{39}(7), 783--797.
https://doi.org/10.1080/0144929X.2019.1623368

\textbf{Journal article with DOI:} Include the DOI as a live hyperlink.
Always include DOI when available.

\textbf{Book:} \textgreater{} Blumler, J. G., \& Katz, E. (1974).
\emph{The uses of mass communications: Current perspectives on
gratifications research}. Sage.

\textbf{Book chapter:} \textgreater{} McQuail, D. (2010). The future of
communication theory. In J. Gripsrud (Ed.), \emph{Relocating television:
Television in the digital context} (pp.~15--28). Routledge.

\textbf{Website / online resource:} \textgreater{} VU Amsterdam. (2024).
\emph{Thesis submission guidelines}.
https://vu.nl/en/student/graduation-and-diploma

\begin{tcolorbox}[enhanced jigsaw, bottomrule=.15mm, opacityback=0, toprule=.15mm, colframe=quarto-callout-warning-color-frame, left=2mm, rightrule=.15mm, title=\textcolor{quarto-callout-warning-color}{\faExclamationTriangle}\hspace{0.5em}{Reference managers often produce incomplete book chapter references}, colbacktitle=quarto-callout-warning-color!10!white, bottomtitle=1mm, colback=white, breakable, toptitle=1mm, titlerule=0mm, leftrule=.75mm, arc=.35mm, opacitybacktitle=0.6, coltitle=black]

Always verify book chapter references manually --- check that editor
names, book title, page numbers, and publisher are all included.

\end{tcolorbox}

\section{Quality of references}\label{quality-of-references}

\begin{itemize}
\tightlist
\item
  The \textbf{large majority} of your references should be peer-reviewed
  journal articles
\item
  Prefer articles from journals with a strong reputation in social
  sciences
\item
  Avoid citing: Wikipedia, blogs, news articles, BA/MA theses, or PhD
  dissertations as primary sources for theoretical claims
\item
  When in doubt about a source's quality, discuss it with your
  supervisor
\end{itemize}

\section{Using Zotero}\label{using-zotero}

\href{https://www.zotero.org}{Zotero} generates APA 7th references
automatically. However, always \textbf{check the output} --- automatic
generation sometimes produces errors or incomplete entries, especially
for: - Book chapters (often missing editors, page numbers, or publisher)
- Conference papers - Preprints

\section{Checklist: Reference list}\label{checklist-reference-list}

\begin{itemize}
\tightlist
\item[$\square$]
  Every in-text citation has a corresponding entry in the reference
  list, and vice versa
\item[$\square$]
  All entries follow APA 7th format
\item[$\square$]
  Journal article titles are in sentence case (capitalize only first
  word, proper nouns, and first word after a colon)
\item[$\square$]
  Journal names are in title case and italicized
\item[$\square$]
  DOIs are included where available
\item[$\square$]
  Book chapter references are complete (editors, page numbers,
  publisher)
\item[$\square$]
  References are in alphabetical order by first author's last name
\end{itemize}

\chapter{Open Science}\label{open-science}

Open science refers to a set of practices that make scientific research
more transparent, reproducible, and accessible. Adopting open science
practices in your thesis strengthens the quality of your work and aligns
with the standards increasingly expected in Communication Science.

\section{Pre-registration}\label{pre-registration-1}

\textbf{Pre-registration} means publicly documenting your research
design, hypotheses, and analysis plan \emph{before} you collect data.
It:

\begin{itemize}
\tightlist
\item
  Prevents \emph{p}-hacking and HARKing (Hypothesizing After Results are
  Known)
\item
  Increases credibility of your confirmatory findings
\item
  Distinguishes what you planned from exploratory analyses you add later
\end{itemize}

\subsection{Where to pre-register}\label{where-to-pre-register}

The standard platform is the \textbf{Open Science Framework (OSF)}:
\href{https://osf.io}{osf.io}

To pre-register: 1. Create a free OSF account (log in with your
institutional email) 2. Create a new project for your thesis 3. Under
\emph{Registrations}, choose a registration template (e.g., AsPredicted
or OSF Prereg) 4. Complete your hypotheses, design, and analysis plan 5.
Submit --- you receive a time-stamped, publicly accessible registration

\begin{tcolorbox}[enhanced jigsaw, bottomrule=.15mm, opacityback=0, toprule=.15mm, colframe=quarto-callout-tip-color-frame, left=2mm, rightrule=.15mm, title=\textcolor{quarto-callout-tip-color}{\faLightbulb}\hspace{0.5em}{Discuss pre-registration with your supervisor}, colbacktitle=quarto-callout-tip-color!10!white, bottomtitle=1mm, colback=white, breakable, toptitle=1mm, titlerule=0mm, leftrule=.75mm, arc=.35mm, opacitybacktitle=0.6, coltitle=black]

Pre-registration is optional for MA theses but encouraged. Your
supervisor can help you decide whether and how to do it. If you
register, include the OSF link in your method section.

\end{tcolorbox}

\section{Open data}\label{open-data}

Sharing your data and analysis code allows others to verify, reuse, and
build on your work.

In your VU thesis, you share data with your supervisor via a secure
surfdrive folder (see \href{ethics.qmd}{Ethics \& Data Management}).
After your thesis is graded, you may optionally deposit anonymized data
in a public repository such as:

\begin{itemize}
\tightlist
\item
  \textbf{OSF} (\href{https://osf.io}{osf.io}) --- general purpose
\item
  \textbf{DANS / Data Archiving and Networked Services}
  (\href{https://dans.knaw.nl}{dans.knaw.nl}) --- Dutch research data
  archive
\item
  \textbf{Zenodo} (\href{https://zenodo.org}{zenodo.org}) ---
  CERN-hosted, general purpose
\end{itemize}

\begin{tcolorbox}[enhanced jigsaw, bottomrule=.15mm, opacityback=0, toprule=.15mm, colframe=quarto-callout-warning-color-frame, left=2mm, rightrule=.15mm, title=\textcolor{quarto-callout-warning-color}{\faExclamationTriangle}\hspace{0.5em}{Never share personal data publicly}, colbacktitle=quarto-callout-warning-color!10!white, bottomtitle=1mm, colback=white, breakable, toptitle=1mm, titlerule=0mm, leftrule=.75mm, arc=.35mm, opacitybacktitle=0.6, coltitle=black]

Before depositing any data publicly, ensure all personal and sensitive
information has been removed. When in doubt, consult your supervisor and
the RERC guidelines.

\end{tcolorbox}

\section{Reproducible analysis}\label{reproducible-analysis}

Reproducibility means your results can be re-created exactly from your
raw data and code. Good practices:

\begin{itemize}
\tightlist
\item
  Keep a \textbf{clean, annotated syntax file} (R or SPSS) from the
  start
\item
  Use \textbf{relative file paths} so your code works on other machines
\item
  Document every data transformation and exclusion decision
\item
  Consider using \textbf{R Markdown} or \textbf{Quarto} to combine
  analysis and write-up
\end{itemize}

For R users: resources are available on the Canvas page \emph{R for
Social Scientists}.

\section{Transparency in reporting}\label{transparency-in-reporting}

\begin{itemize}
\tightlist
\item
  Report \textbf{all tested hypotheses}, including non-significant
  findings
\item
  Clearly distinguish \textbf{confirmatory} (pre-registered) from
  \textbf{exploratory} analyses
\item
  If you deviated from your pre-registered plan, explain why
\item
  Report exact \emph{p} values and effect sizes (not just \emph{p}
  \textless{} .05)
\end{itemize}

\section{Open access}\label{open-access}

After passing your thesis, you are required to upload your thesis to the
\textbf{UBVU thesis database} with public access rights (see
\href{assessment.qmd}{Assessment \& Graduation}). This makes your work
openly available for future researchers and students.

\chapter{Using AI Tools}\label{using-ai-tools}

Generative AI tools are increasingly part of the academic workflow. The
Department of Communication Science at VU Amsterdam permits the use of
these tools --- provided you use them \textbf{responsibly and
transparently}.

\begin{tcolorbox}[enhanced jigsaw, bottomrule=.15mm, opacityback=0, toprule=.15mm, colframe=quarto-callout-important-color-frame, left=2mm, rightrule=.15mm, title=\textcolor{quarto-callout-important-color}{\faExclamation}\hspace{0.5em}{The core principle: Human in the Lead}, colbacktitle=quarto-callout-important-color!10!white, bottomtitle=1mm, colback=white, breakable, toptitle=1mm, titlerule=0mm, leftrule=.75mm, arc=.35mm, opacitybacktitle=0.6, coltitle=black]

You are always the one who decides on the content. You are responsible
for the quality and originality of your work. Generative AI cannot do
your thinking for you --- and it should not.

\end{tcolorbox}

\section{What is responsible use?}\label{what-is-responsible-use}

Responsible use means:

\begin{itemize}
\tightlist
\item
  \textbf{You decide on the content} at all times
\item
  \textbf{You are critical of AI output} --- AI is not a source of
  truth; it can produce inaccurate, incomplete, or superficial
  suggestions
\item
  \textbf{You understand what you are doing} --- use AI tools only when
  you understand the topic well enough to evaluate, improve, or reject
  what the tool suggests
\item
  \textbf{You take full responsibility} for the final submitted version
\end{itemize}

Generative AI can be a useful sparring partner for brainstorming,
gaining an overview, structuring your work, or improving text you have
already written. It cannot replace the reading, thinking, and writing
that constitute your learning process.

\section{Academic misconduct and
fraud}\label{academic-misconduct-and-fraud}

\begin{itemize}
\tightlist
\item
  \textbf{Never copy-paste AI-generated text} into your thesis without
  attribution --- this is plagiarism
\item
  Using AI without transparency is \textbf{academic fraud}
\item
  If your supervisor doubts whether you have used AI responsibly, they
  may discuss this with you; persistent doubts go to the Examination
  Board
\end{itemize}

More information:
\href{https://vu.nl/en/student/examinations/academic-integrity}{vu.nl/en/student/examinations/academic-integrity}

VU rules on generative AI:
\href{https://vu.nl/en/student/examinations/generative-ai-your-use-our-expectations}{vu.nl/en/student/examinations/generative-ai-your-use-our-expectations}

\section{VU-licensed vs.~non-licensed
tools}\label{vu-licensed-vs.-non-licensed-tools}

\subsection{VU-licensed tools
(recommended)}\label{vu-licensed-tools-recommended}

Log in with your \textbf{VU email account} --- your data will not be
used to train the model:

\begin{longtable}[]{@{}
  >{\raggedright\arraybackslash}p{(\columnwidth - 2\tabcolsep) * \real{0.4286}}
  >{\raggedright\arraybackslash}p{(\columnwidth - 2\tabcolsep) * \real{0.5714}}@{}}
\toprule\noalign{}
\begin{minipage}[b]{\linewidth}\raggedright
Tool
\end{minipage} & \begin{minipage}[b]{\linewidth}\raggedright
Access
\end{minipage} \\
\midrule\noalign{}
\endhead
\bottomrule\noalign{}
\endlastfoot
\textbf{Microsoft Copilot} &
\href{https://copilot.microsoft.com}{copilot.microsoft.com} --- log in
with VU credentials \\
\textbf{Google Gemini} &
\href{https://gemini.google.com}{gemini.google.com} --- log in with VU
email \\
\textbf{NotebookLM} &
\href{https://notebooklm.google.com}{notebooklm.google.com} --- log in
with VU email \\
\end{longtable}

\subsection{Non-licensed tools (use at your own
risk)}\label{non-licensed-tools-use-at-your-own-risk}

Tools like \textbf{ChatGPT}, \textbf{Elicit}, \textbf{QuillBot}, and
others fall \textbf{outside VU's license}. When using these: - Your data
may be used for model training - You are personally responsible for
privacy and data protection - \textbf{Never upload personal research
data} (names, email addresses, interview transcripts, survey data) -
\textbf{Never upload paywalled articles} --- this violates copyright

\begin{tcolorbox}[enhanced jigsaw, bottomrule=.15mm, opacityback=0, toprule=.15mm, colframe=quarto-callout-warning-color-frame, left=2mm, rightrule=.15mm, title=\textcolor{quarto-callout-warning-color}{\faExclamationTriangle}\hspace{0.5em}{GDPR and data protection}, colbacktitle=quarto-callout-warning-color!10!white, bottomtitle=1mm, colback=white, breakable, toptitle=1mm, titlerule=0mm, leftrule=.75mm, arc=.35mm, opacitybacktitle=0.6, coltitle=black]

Never upload sensitive or personal data to any AI tool --- not even
VU-licensed ones. This includes: names, email addresses, research data,
financial records, or company details from internships or partnerships.

\end{tcolorbox}

\section{Responsible AI use by thesis
phase}\label{responsible-ai-use-by-thesis-phase}

\begin{longtable}[]{@{}
  >{\raggedright\arraybackslash}p{(\columnwidth - 4\tabcolsep) * \real{0.3333}}
  >{\raggedright\arraybackslash}p{(\columnwidth - 4\tabcolsep) * \real{0.3333}}
  >{\raggedright\arraybackslash}p{(\columnwidth - 4\tabcolsep) * \real{0.3333}}@{}}
\toprule\noalign{}
\begin{minipage}[b]{\linewidth}\raggedright
Thesis Phase
\end{minipage} & \begin{minipage}[b]{\linewidth}\raggedright
Responsible Use
\end{minipage} & \begin{minipage}[b]{\linewidth}\raggedright
Useful Tools
\end{minipage} \\
\midrule\noalign{}
\endhead
\bottomrule\noalign{}
\endlastfoot
\textbf{Topic selection} & Brainstorming ideas, identifying research
gaps & Gemini, Elicit \\
\textbf{Literature review} & Finding keywords/synonyms, searching
articles, generating summaries (always verify originals) & Elicit,
Scite, Research Rabbit, NotebookLM, Undermind \\
\textbf{Research design} & Generating ideas, creating experimental
stimuli (justify all choices yourself) & Gemini, Qualtrics AI \\
\textbf{Proposal writing} & Feedback on text you wrote; rephrasing,
shortening, grammar & Gemini, Grammarly, Wordtune \\
\textbf{Data collection} & In some cases, automating data collection
(verify quality yourself) & Code Interpreter, Google Forms AI \\
\textbf{Data analysis} & Generating/debugging R code, creating
visualizations (verify all output) & Gemini, NVivo, Orange3 \\
\textbf{Writing the thesis} & Improving clarity and structure of your
own text & Gemini, Grammarly, QuillBot, Wordtune \\
\textbf{Revision \& feedback} & Improving academic style, tone,
consistency & Gemini, Grammarly, Hemingway Editor, Trinka AI \\
\end{longtable}

\begin{tcolorbox}[enhanced jigsaw, bottomrule=.15mm, opacityback=0, toprule=.15mm, colframe=quarto-callout-tip-color-frame, left=2mm, rightrule=.15mm, title=\textcolor{quarto-callout-tip-color}{\faLightbulb}\hspace{0.5em}{Sustainability}, colbacktitle=quarto-callout-tip-color!10!white, bottomtitle=1mm, colback=white, breakable, toptitle=1mm, titlerule=0mm, leftrule=.75mm, arc=.35mm, opacitybacktitle=0.6, coltitle=black]

Every query to a generative AI system requires significant computing
energy. Use these tools mindfully and purposefully.

\end{tcolorbox}

\section{The AI statement (required)}\label{the-ai-statement-required}

All theses must include a \textbf{Statement about Generative AI and
AI-Assisted Technology usage} at the end of the thesis, after the
reference list. This statement is \textbf{not counted in the word
count}.

Your statement must answer:

\begin{enumerate}
\def\labelenumi{\arabic{enumi}.}
\tightlist
\item
  \textbf{Did you use generative AI?} (yes/no --- if no, the statement
  ends here)
\item
  \textbf{Which tools did you use, and how?} --- for each tool, specify
  the purpose (e.g., \emph{``I used Google Gemini to receive feedback on
  paragraph structure in the introduction and theory sections''})
\item
  \textbf{How did you handle the AI-generated content?} --- end with a
  statement taking responsibility for the final work
\end{enumerate}

\subsection{Example AI statements}\label{example-ai-statements}

\begin{quote}
\emph{``During the process of writing this thesis, I used Undermind to
explore the scientific literature. I always verified the information
received via AI searches by reading the original sources. I also used
Google Gemini to receive feedback on my own writing regarding paragraph
structure and academic phrasing. The final formulations were written and
checked by myself. I take full responsibility for the final submitted
version.''}
\end{quote}

\begin{quote}
\emph{``When writing the theory sections, I used Grammarly to improve
readability and language. I also used Copilot to debug errors in my R
code during analyses. When using these tools, I carefully reviewed and
edited the generated content. I take full responsibility for the content
of this thesis.''}
\end{quote}

\subsection{In-chapter disclosure}\label{in-chapter-disclosure}

If you used AI for specific chapter-level tasks, also note this within
the relevant chapter. For example: - \emph{Theory section}: used AI to
generate an initial reading list on a topic - \emph{Method section}:
used AI to generate experimental stimuli - \emph{Results section}: used
AI to generate or refine analysis code

Keep a record of your AI prompts and outputs --- you may be asked to
show these.

\chapter{R for Data Analysis}\label{r-for-data-analysis}

R is the primary statistical computing environment used in Communication
Science research at VU Amsterdam. This chapter points you to resources
for learning R, running common analyses, and reporting results in APA
format. It also covers how AI tools can help you write and debug R code.

\section{Getting started with R}\label{getting-started-with-r}

If you are new to R or need a refresher, start here:

\begin{longtable}[]{@{}
  >{\raggedright\arraybackslash}p{(\columnwidth - 2\tabcolsep) * \real{0.4000}}
  >{\raggedright\arraybackslash}p{(\columnwidth - 2\tabcolsep) * \real{0.6000}}@{}}
\toprule\noalign{}
\begin{minipage}[b]{\linewidth}\raggedright
Resource
\end{minipage} & \begin{minipage}[b]{\linewidth}\raggedright
What it covers
\end{minipage} \\
\midrule\noalign{}
\endhead
\bottomrule\noalign{}
\endlastfoot
\textbf{R for Social Scientists} (Canvas) & The departmental resource
--- covers common analyses relevant to Communication Science; start
here \\
\href{https://github.com/ccs-amsterdam/r-course-material}{CCS Amsterdam
R Course Material} & Modular R tutorials developed at VU Amsterdam for
Communication Science students --- covers R fundamentals, tidyverse,
ggplot2, regression, multilevel models, factor analysis, text analysis,
topic modeling, and more \\
\href{https://r4ds.hadley.nz}{R for Data Science} (Wickham et al.) &
Free online book; covers data import, transformation, visualization with
the tidyverse \\
\href{https://learningstatisticswithr.com}{Learning Statistics with R}
(Navarro) & Free textbook; covers statistics concepts alongside R
implementation \\
\href{https://swirlstats.com}{swirl} & Interactive R tutorials you run
inside R itself \\
\href{https://posit.co/resources/cheatsheets/}{Posit Cheatsheets} &
One-page visual references for ggplot2, dplyr, tidyr, etc. \\
\end{longtable}

\begin{tcolorbox}[enhanced jigsaw, bottomrule=.15mm, opacityback=0, toprule=.15mm, colframe=quarto-callout-tip-color-frame, left=2mm, rightrule=.15mm, title=\textcolor{quarto-callout-tip-color}{\faLightbulb}\hspace{0.5em}{Statistics helpdesk}, colbacktitle=quarto-callout-tip-color!10!white, bottomtitle=1mm, colback=white, breakable, toptitle=1mm, titlerule=0mm, leftrule=.75mm, arc=.35mm, opacitybacktitle=0.6, coltitle=black]

VU Amsterdam runs a \textbf{statistics helpdesk} for students. Check
Canvas for the current schedule and how to book an appointment. This is
the right place for questions about which test to use and how to run it
--- not your supervisor.

\end{tcolorbox}

\section{Essential packages}\label{essential-packages}

Install these packages to cover most analyses you'll need:

\begin{Shaded}
\begin{Highlighting}[]
\FunctionTok{install.packages}\NormalTok{(}\FunctionTok{c}\NormalTok{(}
  \StringTok{"tidyverse"}\NormalTok{,   }\CommentTok{\# data manipulation and visualization (dplyr, ggplot2, tidyr, ...)}
  \StringTok{"haven"}\NormalTok{,       }\CommentTok{\# import SPSS (.sav) and Stata files}
  \StringTok{"psych"}\NormalTok{,       }\CommentTok{\# reliability (alpha), descriptive stats, factor analysis}
  \StringTok{"car"}\NormalTok{,         }\CommentTok{\# ANOVA, Levene\textquotesingle{}s test, recoding}
  \StringTok{"lme4"}\NormalTok{,        }\CommentTok{\# multilevel / mixed models}
  \StringTok{"lavaan"}\NormalTok{,      }\CommentTok{\# SEM and mediation/moderation (alternative to PROCESS)}
  \StringTok{"mediation"}\NormalTok{,   }\CommentTok{\# causal mediation analysis}
  \StringTok{"effectsize"}\NormalTok{,  }\CommentTok{\# Cohen\textquotesingle{}s d, eta{-}squared, omega{-}squared}
  \StringTok{"apaTables"}\NormalTok{,   }\CommentTok{\# export APA{-}formatted tables to Word}
  \StringTok{"papaja"}\NormalTok{,      }\CommentTok{\# APA{-}style reporting in R Markdown / Quarto}
  \StringTok{"report"}       \CommentTok{\# automatic plain{-}language results reporting}
\NormalTok{))}
\end{Highlighting}
\end{Shaded}

\section{Common analyses}\label{common-analyses}

\subsection{Descriptive statistics and
correlations}\label{descriptive-statistics-and-correlations}

\begin{Shaded}
\begin{Highlighting}[]
\FunctionTok{library}\NormalTok{(tidyverse)}
\FunctionTok{library}\NormalTok{(psych)}

\CommentTok{\# Descriptive statistics}
\FunctionTok{describe}\NormalTok{(data[, }\FunctionTok{c}\NormalTok{(}\StringTok{"var1"}\NormalTok{, }\StringTok{"var2"}\NormalTok{, }\StringTok{"var3"}\NormalTok{)])}

\CommentTok{\# Correlation matrix (APA{-}ready)}
\FunctionTok{library}\NormalTok{(apaTables)}
\FunctionTok{apa.cor.table}\NormalTok{(data[, }\FunctionTok{c}\NormalTok{(}\StringTok{"var1"}\NormalTok{, }\StringTok{"var2"}\NormalTok{, }\StringTok{"var3"}\NormalTok{)],}
              \AttributeTok{filename =} \StringTok{"Table1\_correlations.doc"}\NormalTok{)}
\end{Highlighting}
\end{Shaded}

\subsection{Reliability (Cronbach's
α)}\label{reliability-cronbachs-ux3b1}

\begin{Shaded}
\begin{Highlighting}[]
\FunctionTok{library}\NormalTok{(psych)}

\CommentTok{\# Select items belonging to one scale}
\FunctionTok{alpha}\NormalTok{(data[, }\FunctionTok{c}\NormalTok{(}\StringTok{"item1"}\NormalTok{, }\StringTok{"item2"}\NormalTok{, }\StringTok{"item3"}\NormalTok{, }\StringTok{"item4"}\NormalTok{)])}
\CommentTok{\# Look at: raw\_alpha (Cronbach\textquotesingle{}s α)}
\end{Highlighting}
\end{Shaded}

\subsection{\texorpdfstring{Independent samples
\emph{t}-test}{Independent samples t-test}}\label{independent-samples-t-test}

\begin{Shaded}
\begin{Highlighting}[]
\CommentTok{\# Test whether group means differ}
\FunctionTok{t.test}\NormalTok{(outcome }\SpecialCharTok{\textasciitilde{}}\NormalTok{ group, }\AttributeTok{data =}\NormalTok{ data, }\AttributeTok{var.equal =} \ConstantTok{TRUE}\NormalTok{)}

\CommentTok{\# Effect size}
\FunctionTok{library}\NormalTok{(effectsize)}
\FunctionTok{cohens\_d}\NormalTok{(outcome }\SpecialCharTok{\textasciitilde{}}\NormalTok{ group, }\AttributeTok{data =}\NormalTok{ data)}
\end{Highlighting}
\end{Shaded}

Report as: \emph{t}(df) = X.XX, \emph{p} = .XXX, \emph{d} = X.XX

\subsection{One-way and factorial
ANOVA}\label{one-way-and-factorial-anova}

\begin{Shaded}
\begin{Highlighting}[]
\FunctionTok{library}\NormalTok{(car)}

\CommentTok{\# Factorial ANOVA (e.g., 2x2 between{-}subjects experiment)}
\NormalTok{model }\OtherTok{\textless{}{-}} \FunctionTok{lm}\NormalTok{(outcome }\SpecialCharTok{\textasciitilde{}}\NormalTok{ factor\_a }\SpecialCharTok{*}\NormalTok{ factor\_b, }\AttributeTok{data =}\NormalTok{ data)}
\FunctionTok{Anova}\NormalTok{(model, }\AttributeTok{type =} \StringTok{"III"}\NormalTok{)   }\CommentTok{\# Type III sums of squares}

\CommentTok{\# Effect size}
\FunctionTok{library}\NormalTok{(effectsize)}
\FunctionTok{eta\_squared}\NormalTok{(model, }\AttributeTok{partial =} \ConstantTok{TRUE}\NormalTok{)}
\end{Highlighting}
\end{Shaded}

Report as: \emph{F}(df1, df2) = X.XX, \emph{p} = .XXX, η²\_p = .XX

\subsection{Linear regression}\label{linear-regression}

\begin{Shaded}
\begin{Highlighting}[]
\CommentTok{\# Simple regression}
\NormalTok{model1 }\OtherTok{\textless{}{-}} \FunctionTok{lm}\NormalTok{(outcome }\SpecialCharTok{\textasciitilde{}}\NormalTok{ predictor, }\AttributeTok{data =}\NormalTok{ data)}
\FunctionTok{summary}\NormalTok{(model1)}

\CommentTok{\# Multiple regression}
\NormalTok{model2 }\OtherTok{\textless{}{-}} \FunctionTok{lm}\NormalTok{(outcome }\SpecialCharTok{\textasciitilde{}}\NormalTok{ predictor1 }\SpecialCharTok{+}\NormalTok{ predictor2 }\SpecialCharTok{+}\NormalTok{ covariate, }\AttributeTok{data =}\NormalTok{ data)}
\FunctionTok{summary}\NormalTok{(model2)}

\CommentTok{\# Standardized coefficients}
\FunctionTok{library}\NormalTok{(effectsize)}
\FunctionTok{standardize\_parameters}\NormalTok{(model2)}
\end{Highlighting}
\end{Shaded}

Report as: \emph{b} = X.XX, \emph{SE} = X.XX, \emph{β} = X.XX,
\emph{t}(df) = X.XX, \emph{p} = .XXX

\subsection{Moderation (interaction
effects)}\label{moderation-interaction-effects}

\begin{Shaded}
\begin{Highlighting}[]
\CommentTok{\# Center continuous predictors before creating interaction term}
\NormalTok{data}\SpecialCharTok{$}\NormalTok{pred\_c     }\OtherTok{\textless{}{-}} \FunctionTok{scale}\NormalTok{(data}\SpecialCharTok{$}\NormalTok{predictor, }\AttributeTok{center =} \ConstantTok{TRUE}\NormalTok{, }\AttributeTok{scale =} \ConstantTok{FALSE}\NormalTok{)}
\NormalTok{data}\SpecialCharTok{$}\NormalTok{moderator\_c }\OtherTok{\textless{}{-}} \FunctionTok{scale}\NormalTok{(data}\SpecialCharTok{$}\NormalTok{moderator, }\AttributeTok{center =} \ConstantTok{TRUE}\NormalTok{, }\AttributeTok{scale =} \ConstantTok{FALSE}\NormalTok{)}

\NormalTok{model\_mod }\OtherTok{\textless{}{-}} \FunctionTok{lm}\NormalTok{(outcome }\SpecialCharTok{\textasciitilde{}}\NormalTok{ pred\_c }\SpecialCharTok{*}\NormalTok{ moderator\_c, }\AttributeTok{data =}\NormalTok{ data)}
\FunctionTok{summary}\NormalTok{(model\_mod)}

\CommentTok{\# Visualize interaction with ggplot2 or the \textquotesingle{}interactions\textquotesingle{} package}
\FunctionTok{library}\NormalTok{(interactions)}
\FunctionTok{interact\_plot}\NormalTok{(model\_mod, }\AttributeTok{pred =}\NormalTok{ pred\_c, }\AttributeTok{modx =}\NormalTok{ moderator\_c)}
\end{Highlighting}
\end{Shaded}

\subsection{Mediation analysis (PROCESS /
lavaan)}\label{mediation-analysis-process-lavaan}

\begin{Shaded}
\begin{Highlighting}[]
\CommentTok{\# Using the \textquotesingle{}mediation\textquotesingle{} package (bootstrapped indirect effects)}
\FunctionTok{library}\NormalTok{(mediation)}

\CommentTok{\# Step 1: model for mediator}
\NormalTok{m\_model }\OtherTok{\textless{}{-}} \FunctionTok{lm}\NormalTok{(mediator }\SpecialCharTok{\textasciitilde{}}\NormalTok{ predictor, }\AttributeTok{data =}\NormalTok{ data)}

\CommentTok{\# Step 2: model for outcome}
\NormalTok{y\_model }\OtherTok{\textless{}{-}} \FunctionTok{lm}\NormalTok{(outcome }\SpecialCharTok{\textasciitilde{}}\NormalTok{ predictor }\SpecialCharTok{+}\NormalTok{ mediator, }\AttributeTok{data =}\NormalTok{ data)}

\CommentTok{\# Bootstrap mediation}
\NormalTok{med\_out }\OtherTok{\textless{}{-}} \FunctionTok{mediate}\NormalTok{(m\_model, y\_model,}
                   \AttributeTok{treat =} \StringTok{"predictor"}\NormalTok{, }\AttributeTok{mediator =} \StringTok{"mediator"}\NormalTok{,}
                   \AttributeTok{boot =} \ConstantTok{TRUE}\NormalTok{, }\AttributeTok{sims =} \DecValTok{5000}\NormalTok{)}
\FunctionTok{summary}\NormalTok{(med\_out)}
\end{Highlighting}
\end{Shaded}

Alternatively, use the \textbf{PROCESS macro} (Hayes) in SPSS, or
\textbf{lavaan} for structural equation modeling in R.

\subsection{Manipuation check
(experiment)}\label{manipuation-check-experiment}

\begin{Shaded}
\begin{Highlighting}[]
\CommentTok{\# Did participants in different conditions perceive the manipulation differently?}
\FunctionTok{t.test}\NormalTok{(manipulation\_check }\SpecialCharTok{\textasciitilde{}}\NormalTok{ condition, }\AttributeTok{data =}\NormalTok{ data)}

\CommentTok{\# For 2x2 design}
\NormalTok{aov\_result }\OtherTok{\textless{}{-}} \FunctionTok{aov}\NormalTok{(manipulation\_check }\SpecialCharTok{\textasciitilde{}}\NormalTok{ factor\_a }\SpecialCharTok{*}\NormalTok{ factor\_b, }\AttributeTok{data =}\NormalTok{ data)}
\FunctionTok{summary}\NormalTok{(aov\_result)}
\end{Highlighting}
\end{Shaded}

\section{APA-formatted reporting in
R}\label{apa-formatted-reporting-in-r}

\subsection{\texorpdfstring{The \texttt{report}
package}{The report package}}\label{the-report-package}

\begin{Shaded}
\begin{Highlighting}[]
\FunctionTok{library}\NormalTok{(report)}

\NormalTok{model }\OtherTok{\textless{}{-}} \FunctionTok{lm}\NormalTok{(outcome }\SpecialCharTok{\textasciitilde{}}\NormalTok{ predictor, }\AttributeTok{data =}\NormalTok{ data)}
\FunctionTok{report}\NormalTok{(model)}
\CommentTok{\# Returns a plain{-}language sentence you can adapt for your results section}
\end{Highlighting}
\end{Shaded}

\subsection{\texorpdfstring{The \texttt{apaTables}
package}{The apaTables package}}\label{the-apatables-package}

\begin{Shaded}
\begin{Highlighting}[]
\FunctionTok{library}\NormalTok{(apaTables)}

\CommentTok{\# APA regression table exported to Word}
\FunctionTok{apa.reg.table}\NormalTok{(model1, model2,}
              \AttributeTok{filename =} \StringTok{"Table2\_regression.doc"}\NormalTok{)}
\end{Highlighting}
\end{Shaded}

\subsection{papaja (APA manuscripts in Quarto/R
Markdown)}\label{papaja-apa-manuscripts-in-quartor-markdown}

If you are writing your thesis in R Markdown or Quarto, the
\href{https://frederikaust.com/papaja_man/}{papaja} package provides
APA-formatted output. Use \texttt{apa\_print()} to inline statistics:

\begin{Shaded}
\begin{Highlighting}[]
\FunctionTok{library}\NormalTok{(papaja)}
\NormalTok{result }\OtherTok{\textless{}{-}} \FunctionTok{t.test}\NormalTok{(outcome }\SpecialCharTok{\textasciitilde{}}\NormalTok{ group, }\AttributeTok{data =}\NormalTok{ data)}
\FunctionTok{apa\_print}\NormalTok{(result)}\SpecialCharTok{$}\NormalTok{statistic   }\CommentTok{\# → "t(48) = 2.34, p = .023"}
\end{Highlighting}
\end{Shaded}

\section{Keeping your syntax clean}\label{keeping-your-syntax-clean}

Good syntax habits matter --- your supervisor needs to be able to
reproduce your analyses from your raw data.

\begin{Shaded}
\begin{Highlighting}[]
\CommentTok{\# ============================================================}
\CommentTok{\# MA Thesis Analysis — [Your Name] — [Date]}
\CommentTok{\# ============================================================}

\CommentTok{\# 0. Load packages {-}{-}{-}{-}}
\FunctionTok{library}\NormalTok{(tidyverse)}
\FunctionTok{library}\NormalTok{(psych)}
\FunctionTok{library}\NormalTok{(effectsize)}

\CommentTok{\# 1. Import data {-}{-}{-}{-}}
\NormalTok{data\_raw }\OtherTok{\textless{}{-}}\NormalTok{ haven}\SpecialCharTok{::}\FunctionTok{read\_sav}\NormalTok{(}\StringTok{"data/raw\_data.sav"}\NormalTok{)}

\CommentTok{\# 2. Data cleaning {-}{-}{-}{-}}
\CommentTok{\# Remove incomplete responses (n = 9)}
\NormalTok{data\_clean }\OtherTok{\textless{}{-}}\NormalTok{ data\_raw }\SpecialCharTok{|\textgreater{}}
  \FunctionTok{filter}\NormalTok{(progress }\SpecialCharTok{==} \DecValTok{100}\NormalTok{)}

\CommentTok{\# 3. Compute scale scores {-}{-}{-}{-}}
\CommentTok{\# Scale: Vaccine Hesitancy (α = .82 in original; check below)}
\NormalTok{data\_clean }\OtherTok{\textless{}{-}}\NormalTok{ data\_clean }\SpecialCharTok{|\textgreater{}}
  \FunctionTok{mutate}\NormalTok{(}\AttributeTok{vacc\_hes =} \FunctionTok{rowMeans}\NormalTok{(}\FunctionTok{across}\NormalTok{(}\FunctionTok{c}\NormalTok{(vh1, vh2, vh3, vh4, vh5)),}
                              \AttributeTok{na.rm =} \ConstantTok{TRUE}\NormalTok{))}

\CommentTok{\# 4. Descriptive statistics {-}{-}{-}{-}}
\FunctionTok{describe}\NormalTok{(data\_clean[, }\FunctionTok{c}\NormalTok{(}\StringTok{"vacc\_hes"}\NormalTok{, }\StringTok{"soc\_media\_exp"}\NormalTok{, }\StringTok{"age"}\NormalTok{)])}

\CommentTok{\# 5. H1: Regression of social media exposure on vaccine hesitancy {-}{-}{-}{-}}
\NormalTok{h1\_model }\OtherTok{\textless{}{-}} \FunctionTok{lm}\NormalTok{(vacc\_hes }\SpecialCharTok{\textasciitilde{}}\NormalTok{ soc\_media\_exp, }\AttributeTok{data =}\NormalTok{ data\_clean)}
\FunctionTok{summary}\NormalTok{(h1\_model)}
\end{Highlighting}
\end{Shaded}

\begin{tcolorbox}[enhanced jigsaw, bottomrule=.15mm, opacityback=0, toprule=.15mm, colframe=quarto-callout-tip-color-frame, left=2mm, rightrule=.15mm, title=\textcolor{quarto-callout-tip-color}{\faLightbulb}\hspace{0.5em}{Use section headers in your script}, colbacktitle=quarto-callout-tip-color!10!white, bottomtitle=1mm, colback=white, breakable, toptitle=1mm, titlerule=0mm, leftrule=.75mm, arc=.35mm, opacitybacktitle=0.6, coltitle=black]

The \texttt{\#\ Section\ -\/-\/-\/-} pattern (four dashes) creates
collapsible sections in RStudio's code outline. This keeps your syntax
file navigable when it grows long.

\end{tcolorbox}

\section{Getting help with R code}\label{getting-help-with-r-code}

Beyond the statistics helpdesk and Canvas resources, AI tools can be
genuinely useful for R:

\begin{longtable}[]{@{}
  >{\raggedright\arraybackslash}p{(\columnwidth - 2\tabcolsep) * \real{0.2727}}
  >{\raggedright\arraybackslash}p{(\columnwidth - 2\tabcolsep) * \real{0.7273}}@{}}
\toprule\noalign{}
\begin{minipage}[b]{\linewidth}\raggedright
Task
\end{minipage} & \begin{minipage}[b]{\linewidth}\raggedright
Useful approach
\end{minipage} \\
\midrule\noalign{}
\endhead
\bottomrule\noalign{}
\endlastfoot
Debugging an error message & Paste the error + relevant code into Gemini
or Copilot and ask what is wrong \\
Finding the right function & Describe what you want to do in plain
English \\
Writing boilerplate (e.g., a for-loop) & Ask AI to generate it, then
check you understand it \\
Interpreting output & Ask AI to explain what a specific output means \\
\end{longtable}

\begin{tcolorbox}[enhanced jigsaw, bottomrule=.15mm, opacityback=0, toprule=.15mm, colframe=quarto-callout-warning-color-frame, left=2mm, rightrule=.15mm, title=\textcolor{quarto-callout-warning-color}{\faExclamationTriangle}\hspace{0.5em}{Verify AI-generated R code}, colbacktitle=quarto-callout-warning-color!10!white, bottomtitle=1mm, colback=white, breakable, toptitle=1mm, titlerule=0mm, leftrule=.75mm, arc=.35mm, opacitybacktitle=0.6, coltitle=black]

AI tools can produce plausible-looking but incorrect R code. Always run
it, check the output makes sense, and make sure you understand what
every line does before including results in your thesis. See
\href{ai-tools.qmd}{Using AI Tools} for the full policy.

\end{tcolorbox}

\part{Submission \& Assessment}

\chapter{Assessment \& Graduation}\label{assessment-graduation}

\section{Submitting your final
thesis}\label{submitting-your-final-thesis}

Submit your final thesis \textbf{no later than the first-opportunity
deadline} (see \href{timeline.qmd}{Timeline \& Deadlines}):

\begin{enumerate}
\def\labelenumi{\arabic{enumi}.}
\tightlist
\item
  \textbf{Upload on Canvas} (required --- the thesis is checked for
  plagiarism here)
\item
  \textbf{Email to your supervisor and second reviewer} (confirm with
  your supervisor whether a printed copy is also required)
\item
  \textbf{Ensure your supervisor has access} to your final data file and
  syntax via the shared surfdrive folder
\end{enumerate}

\section{How your thesis is assessed}\label{how-your-thesis-is-assessed}

Your thesis is evaluated by \textbf{two reviewers}: your supervisor
(first reviewer) and a second reviewer. Both use the official
\textbf{assessment form} (available on Canvas).

\begin{longtable}[]{@{}
  >{\raggedright\arraybackslash}p{(\columnwidth - 2\tabcolsep) * \real{0.5000}}
  >{\raggedright\arraybackslash}p{(\columnwidth - 2\tabcolsep) * \real{0.5000}}@{}}
\toprule\noalign{}
\begin{minipage}[b]{\linewidth}\raggedright
\end{minipage} & \begin{minipage}[b]{\linewidth}\raggedright
Details
\end{minipage} \\
\midrule\noalign{}
\endhead
\bottomrule\noalign{}
\endlastfoot
\textbf{Time to grade} & Two weeks (10 working days) after submission
deadline \\
\textbf{Final grade} & Average of first and second reviewer \\
\textbf{Disagreement} & If assessors differ by more than 2 points, or
one fails and the other passes → third assessor assigned by the
Examination Board (binding) \\
\textbf{Passing grade} & ≥ 5.5 \\
\end{longtable}

\section{Resits and late submission}\label{resits-and-late-submission}

Failing to submit at the first deadline is treated the same as missing
an exam --- you have used one exam opportunity.

\begin{longtable}[]{@{}
  >{\raggedright\arraybackslash}p{(\columnwidth - 2\tabcolsep) * \real{0.5000}}
  >{\raggedright\arraybackslash}p{(\columnwidth - 2\tabcolsep) * \real{0.5000}}@{}}
\toprule\noalign{}
\begin{minipage}[b]{\linewidth}\raggedright
Situation
\end{minipage} & \begin{minipage}[b]{\linewidth}\raggedright
What happens
\end{minipage} \\
\midrule\noalign{}
\endhead
\bottomrule\noalign{}
\endlastfoot
\textbf{Not submitted at first deadline} & First exam opportunity used;
one resit available \\
\textbf{Submitted but grade \textless{} 5.5} & One resit available \\
\textbf{Resit submitted but grade \textless{} 5.5} & Failed the course;
may be eligible for a second resit in the next academic year (see
below) \\
\textbf{Not submitted at first deadline AND resit} & Failed the course;
may be eligible for a second resit in the next academic year \\
\end{longtable}

\subsection{Second resit (next academic
year)}\label{second-resit-next-academic-year}

In limited circumstances, you may be eligible for a second resit in the
following academic year (deadline: end of November). This requires
re-enrollment and tuition fee payment. Eligibility is assessed by the
supervisor and thesis coordinator. Contact your study advisor.

\section{Graduation procedure}\label{graduation-procedure}

Once your thesis has a passing grade and you have completed all other
courses:

\begin{enumerate}
\def\labelenumi{\arabic{enumi}.}
\tightlist
\item
  You will automatically receive a graduation email within approximately
  two weeks
\item
  \textbf{Upload your thesis to the UBVU thesis database} --- this is
  mandatory for receiving your diploma

  \begin{itemize}
  \tightlist
  \item
    Choose \emph{Publicly available} under access rights
  \item
    Instructions:
    \href{https://vu.nl/en/about-vu/divisions/university-library/more-about/uploading-your-thesis}{vu.nl/en/about-vu/divisions/university-library/more-about/uploading-your-thesis}
  \end{itemize}
\item
  Indicate your attendance at the \textbf{graduation ceremony} (usually
  in October/November)
\end{enumerate}

\begin{tcolorbox}[enhanced jigsaw, bottomrule=.15mm, opacityback=0, toprule=.15mm, colframe=quarto-callout-note-color-frame, left=2mm, rightrule=.15mm, title=\textcolor{quarto-callout-note-color}{\faInfo}\hspace{0.5em}{No defense required}, colbacktitle=quarto-callout-note-color!10!white, bottomtitle=1mm, colback=white, breakable, toptitle=1mm, titlerule=0mm, leftrule=.75mm, arc=.35mm, opacitybacktitle=0.6, coltitle=black]

There is no oral presentation or defense of the MA thesis in
Communication Science. The diploma ceremony is a celebration, not an
assessment event.

\end{tcolorbox}

\section{Where to find formal
regulations}\label{where-to-find-formal-regulations}

The formal \emph{Regulation Ma thesis} governs all MA theses at the
Faculty of Social Sciences. It is available on VUnet. For the vast
majority of students, these regulations remain in the background ---
supervisor and student find a good working relationship without needing
to consult them. They become relevant if serious disagreements arise.

\section{Getting help}\label{getting-help}

\begin{longtable}[]{@{}
  >{\raggedright\arraybackslash}p{(\columnwidth - 2\tabcolsep) * \real{0.4000}}
  >{\raggedright\arraybackslash}p{(\columnwidth - 2\tabcolsep) * \real{0.6000}}@{}}
\toprule\noalign{}
\begin{minipage}[b]{\linewidth}\raggedright
Need
\end{minipage} & \begin{minipage}[b]{\linewidth}\raggedright
Contact
\end{minipage} \\
\midrule\noalign{}
\endhead
\bottomrule\noalign{}
\endlastfoot
Academic/personal difficulties & Study advisor:
\href{https://vu.nl/en/student/contact-student-guidance-and-support/academic-advisors-fss}{vu.nl/en/student/contact-student-guidance-and-support/academic-advisors-fss} \\
Financial support for delays & Profileringsfonds:
\href{https://vu.nl/en/student/finances/profile-fund}{vu.nl/en/student/finances/profile-fund} \\
Graduation questions &
\href{https://vu.nl/en/student/graduation-and-diploma}{vu.nl/en/student/graduation-and-diploma} \\
Ethics self-check questions & rerc.ssc@vu.nl \\
\end{longtable}


\backmatter

\end{document}
